\documentclass[11pt]{article}
\usepackage[letterpaper,margin=0.83in]{geometry}
\usepackage[T1]{fontenc}
\usepackage{lmodern,amsmath,amssymb,amsthm,mathtools}
\usepackage{booktabs,longtable,array,microtype,xcolor,enumitem,fancyhdr,hyperref}
\hypersetup{colorlinks=true,linkcolor=blue!45!black,urlcolor=blue!55!black,citecolor=blue!45!black,
pdftitle={Affine parameters and a linear CRT selector},
pdfauthor={Research contribution prepared for Vladimir Reshetnikov},
pdfsubject={Explicit quartic equations for syntactic consistency of ZFC}}
\setlength{\parindent}{0pt}
\setlength{\parskip}{5pt}
\setlength{\headheight}{14pt}
\pagestyle{fancy}\fancyhf{}
\fancyhead[L]{\small Affine parameters for Diophantine consistency}
\fancyhead[R]{\small 25 October 2026}\fancyfoot[C]{\thepage}
\newtheorem{theorem}{Theorem}[section]
\newtheorem{lemma}[theorem]{Lemma}
\newtheorem{proposition}[theorem]{Proposition}
\newtheorem{corollary}[theorem]{Corollary}
\theoremstyle{remark}\newtheorem{remark}[theorem]{Remark}
\newcommand{\Z}{\mathbb Z}\newcommand{\N}{\mathbb N}
\newcommand{\Con}{\operatorname{Con}}
\newcommand{\lcm}{\operatorname{lcm}}
\newcommand{\LL}{\mathcal L}
\newcommand{\repo}{\texttt{7a8da414376a80b90ecf799f37ceb26fb9ac674e}}
\newcommand{\PrimaryVars}{238}
\newcommand{\PrimaryRows}{174}
\newcommand{\PrimaryTerms}{5\,479}
\newcommand{\PrimaryHeight}{7\,968}
\newcommand{\PrimaryLastVar}{237}
\newcommand{\PrimaryLastRow}{173}
\newcommand{\SmallVars}{237}
\newcommand{\SmallLastVar}{236}
\newcommand{\SmallRows}{173}
\newcommand{\SmallLastRow}{172}

\title{Affine parameters and a linear CRT selector\\[3pt]
\large Explicit quartic equations for syntactic consistency of ZFC}
\author{Research contribution prepared for Vladimir Reshetnikov\\
\small Developed with AI assistance; ordinary proofs and exact finite certificates}
\date{25 October 2026}
\begin{document}
\maketitle
\begin{abstract}
We simplify the direct numerical proof-calculus construction in the pinned
ProveIt repository. The main polynomial has \PrimaryVars\ unrestricted
integer variables, degree four, \PrimaryRows\ quadratic residuals,
\PrimaryTerms\ collected monomials, and coefficient height \PrimaryHeight.
A companion uses \SmallVars\ integer variables. All coefficients and
residuals are supplied literally, and both principal residual lists are printed here.
The ordinary mathematical proof, composed with the retained C222/ZF/ZFC
calibration, establishes in ZFC that each equation has no integer solution
if and only if $\Con(\mathrm{ZFC})$.

The new mechanism recovers order from already present signed power
certificates. It permits affine parameters for the radix, the beta
multiplier, and the proof horizon. An elementary
least-common-multiple bound also replaces the squared binomial selector
offset by its unsquared growth register. These changes remove witnesses
without introducing a new power request or a new logical encoding.
This is an improvement within the direct small-coefficient quartic family,
not a global minimum, a consistency decision, or a new end-to-end
proof-assistant formalization.
\end{abstract}
\begingroup\small\setcounter{tocdepth}{1}
\setlength{\parskip}{0pt}\tableofcontents
\endgroup
\newpage

\section{The result and its precise meaning}

The requested equivalence concerns \emph{syntactic} consistency. Fix the
repository's recursively presented ZFC axiom predicate and its finite
proof relation. Then $\Con(\mathrm{ZFC})$ means that no finite derivation
of falsity from those axioms exists. We construct ordinary polynomials
with fixed integer coefficients, not polynomials whose coefficients
depend on an undecided consistency question.

For the main profile let $R_0,\ldots,R_{\PrimaryLastRow}$ be the complete
quadratic residuals in Appendix~\ref{app:residuals}. Define
\begin{equation}\label{eq:main}
 P_{\PrimaryVars}(x_0,\ldots,x_{\PrimaryLastVar})
   =\sum_{j=0}^{\PrimaryLastRow}R_j(x_0,\ldots,x_{\PrimaryLastVar})^2.
\end{equation}
Every coordinate ranges over $\Z$. The finite sum introduces no additional
unknown. The endpoint named \texttt{primary} in \texttt{ENDPOINTS.json} identifies the canonical
literal residual list; \texttt{expanded\_polynomial.json} supplies every
collected coefficient of the same polynomial.

\begin{theorem}[Main equivalence]\label{thm:main}
For the polynomial in \eqref{eq:main}, and for every companion in the
endpoint table, the following equivalence has an ordinary proof
formalizable in ZFC:
\[
 \mathrm{ZFC}\vdash
 \left[\neg\exists\boldsymbol x\in\Z^n\;
 P_n(\boldsymbol x)=0\right]\ \longleftrightarrow\ \Con(\mathrm{ZFC}).
\]
No assumption of consistency or inconsistency is used to prove the
equivalence.
\end{theorem}

\subsection{Comparison and the chosen notion of simplicity}

\begin{center}\small
\begin{tabular}{@{}lrrrr@{}}
\toprule Construction & Variables & Rows & Monomials & Height \\
\midrule
Inherited factorial-free & 242 & 178 & 5\,018 & 7\,968 \\
Inherited factorial-free, dense & 241 & 177 & 12\,774 & 7\,968 \\
Incoming shared factorial error & 240 & 177 & 21\,482 & 6\,144 \\
Incoming normalized residues & 244 & 176 & 5\,212 & 3\,968 \\
\midrule
Affine radix and beta & 239 & 175 & 5\,017 & 7\,968 \\
Affine horizon, balanced (main) & 238 & 174 & 5\,479 & 7\,968 \\
Affine horizon, fewer variables & 237 & 173 & 14\,129 & 7\,968 \\
\bottomrule\end{tabular}\end{center}


Height is the largest absolute coefficient \emph{after} collecting the
whole quartic. Residual height is a different quantity. The number of
residuals is not an arithmetic circuit operation count. All comparisons in
the table use unrestricted integer variables and degree four.

The comparison snapshot is ProveIt commit \repo. Its seven incoming ZIPs
were inspected in addition to the integrated collection. In particular,
the incoming Shared Factorial Error construction already had a
240-integer quartic, and the Factorial-Free Binomial CRT construction
already removed factorial authenticity. We credit and retain those
results rather than presenting them as contributions of this report.

The direct family is attractive because its coefficients are inspectable
small integers and its logical source is a numerical proof calculus.
Other repository constructions already have far fewer variables using
different coefficient architectures. The general literature also has
universal reductions with eleven integer unknowns~\cite{bayer}.
Consequently no unrestricted variable-count record or metric-independent
``simplest equation'' is claimed here.

\section{The retained compiler boundary}

The predecessor compiles a finite C222 refutation into bounded arithmetic,
then uses a binomial CRT certificate and shared Pell certificates to
obtain a fixed quadratic system. We retain its local numerical rule
polynomial, all eighteen power requests, all fourteen principal-pair
orientation conditions (with an invertible error-basis change), the shared remainder norm, and the exact integral
projection machinery.

The logical calibration is the composition of the existing interfaces
\begin{align*}
 \texttt{KEndpoint.c222\_iff\_zf\_refutation},\\
 \texttt{ConstructibleChoice.zf\_refutation\_iff\_zfc\_refutation}.
\end{align*}
The first identifies derivability of the falsity code in C222 with a
finite ZF refutation; the second identifies finite ZF and ZFC refutations.
The direct construction uses this calculus itself. It does not encode a
K215 run and does not specialize Jones's universal polynomial.

For clarity, variable exponentiation and binomial notation in the proof
below describe the \emph{semantics} of ordinary polynomial certificates.
Neither operation occurs as a primitive in the delivered quartics.
All requested powers are implemented by the retained Pell graphs.

The principal semantic registers are:
\begin{center}
\begin{tabular}{ll}
\toprule
Register & Meaning after its certificate is authenticated\\
\midrule
$t$ & padded number of local proof rows\\
$F$ & large evaluation point for label multiplicities\\
$y$ & beta modulus scale\\
$x$ & common beta code and amplification exponent\\
$B$ & binomial modulus selecting every proof row\\
$N$ & $B^x$\\
$C,A_1,\ldots,A_8$ & actual CRT representatives\\
$H$ & $2+C+\sum_{i=1}^{8}A_i$\\
$q$ & radix $(BH)^x$\\
\bottomrule
\end{tabular}
\end{center}
Some stored coordinates are translated versions of these registers.
Semantic names alone do not specify that translation; the exact
restoration certificates do.

\section{An unsquared binomial selector}

Let $L_t=\lcm(1,\ldots,t)$. The factorial-free predecessor proves its
CRT selector theorem whenever the offset $R$ satisfies $R\ge L_t^2$.
It used $R=F^2$. Its growth register is already large enough to take
$R=F$ instead.

\begin{lemma}[Elementary lcm estimate]\label{lem:lcm}
For every $n\ge1$, $L_n\le4^{n-1}$.
\end{lemma}
\begin{proof}
Put $m=\lceil n/2\rceil$. For any prime $p$, at most one power of $p$
lies in $(m,n]$. Thus $L_n/L_m$ is squarefree. If $p$ divides this quotient,
its largest power $p^a\le n$ lies in $(m,n]$, and the corresponding
Legendre summand in $v_p\binom nm$ is one. All other summands are
nonnegative. Therefore $L_n/L_m\mid\binom nm$.
For $n\ge2$, each binomial coefficient is at most the sum of coefficients
of the same parity, namely $2^{n-1}$. Starting from $L_1=1$, induction gives
\[
 L_n\le4^{m-1}2^{n-1}\le4^{n-1}.
\]
For even $n$ the first exponent of two is $2n-3$; for odd $n$ it is
$2n-2$. Both give the stated bound.
\end{proof}

This elementary exponential estimate belongs to the classical lcm and
binomial-coefficient toolkit~\cite{farhi,chan}. Our contribution here is
its use to delete the squared selector register.

\begin{corollary}\label{cor:selector}
Each of the following choices satisfies $F>L_t^2$:
\[
\begin{array}{ll}
 F=[2(t+1)]^t,&t\ge8,\\
 F=(t+2)^t,&t\ge16.
\end{array}
\]
\end{corollary}
\begin{proof}
Lemma~\ref{lem:lcm} gives $L_t^2\le16^{t-1}$. In the first case the
base is at least 18, and in the second it is again at least 18.
\end{proof}

We recall exactly the selector property being used. If
$B=\binom{R+t}{t}$ and $R\ge L_t^2$, then every row $i\in\{1,\ldots,t\}$
has a full prime power $p^e\Vert B$ such that $p^e\mid R+i$ and $p^e>t$.
Conversely, every full prime power of $B$ divides at least one $R+i$.
Both directions are necessary: the first ensures soundness for every
row; the second makes the reverse CRT assignment possible at every
prime of $B$.

For completeness, the first assertion follows by choosing
$v=v_p(R+i)>2a$, where $a=\lfloor\log_p t\rfloor$. Such a prime exists
because $R+i>L_t^2$. For $j\ne i$, $v_p(R+j)=v_p(j-i)$, so
\[
 e=v_p(B)=v-d,\qquad
 d=v_p\!\left(t\binom{t-1}{i-1}\right)\le a.
\]
Hence $e>a$ and $e\le v$. The displayed bound on $d$ follows directly
from Legendre's formula: its summands are zero or one and vanish above
$p^a\le t$. For the reverse assertion, let $v=\max_i v_p(R+i)$.
The valuation of $B$ is a sum of zero-one carry terms, all zero above
level $v$; hence $e\le v$, and a maximizing row works.

The new first digit is therefore
\begin{equation}\label{eq:newB}
 B=\binom{F+t}{t},\qquad B\mid C+F+1.
\end{equation}
Both occurrences of the offset change together. The corresponding raw
quadratic row $R-F^2=0$ becomes affine, $R-F=0$, and its coordinate is
eliminated exactly.

\section{Signed powers can recover order}

\subsection{The retained power interface}

Write $X_k(A),Y_k(A)$ for the Pell polynomials, with $X_0=1,Y_0=0$ and
\[
 X_{k+1}=AX_k+(A^2-1)Y_k,\qquad
 Y_{k+1}=X_k+AY_k.
\]
For signed indices use $X_{-k}=X_k$ and $Y_{-k}=-Y_k$. After the
principal pairs have been oriented, a request for $m=n^k$ has modulus
$T_n=2An-n^2-1$. The shared square frame bounds every actual base,
index, and output atom in absolute value by $W$, and gives
$|A|>W^W$.

\begin{lemma}[Nonunit bases force nonnegative exponents]\label{lem:power}
Under the retained oriented Pell interface, if $|n|\ge2$, a requested
integer output forces $k\ge0$ and $m=n^k$, even when $n,A,m,k$ were
initially signed.
\end{lemma}
\begin{proof}
The Pell recurrence gives $m\equiv n^k\pmod{T_n}$ for $k\ge0$.
The absolute difference is smaller than $|T_n|>W^W+W$, so it is zero.
If $k=-j<0$, the inverse-unit recurrence gives
$mn^j\equiv1\pmod{T_n}$. The inverse is legitimate because
$T_n\equiv-1\pmod n$. The absolute error is less than $W^W+1$,
so $mn^j=1$ in $\Z$, impossible when $|n|\ge2$ and $j\ge1$.
\end{proof}

The zero and unit bases are genuine exceptions. The following progressions
are chosen specifically to avoid applying the lemma to them.

\subsection{Affine radix and beta multiplier}

Replace the predecessor's two square parameters by
\begin{equation}\label{eq:affine}
             q=4z+1,\qquad y=F(2v+1),\qquad z,v\in\Z.
\end{equation}
The expression $y$ is now bilinear. The proof cannot initially assume
that either $q$ or $y$ is positive. The key unconditional fact is
$|q+1|=|4z+2|\ge2$.

Retain the two beta endpoints and the self-radix request:
\begin{align}
 x&=F+1+(y+1)u_0,\label{eq:first}\\
 x&=1+[1+(t+1)y]u_t,\label{eq:last}\\
 q&=(BH)^x.\label{eq:radix}
\end{align}
Also retain the existing power request with base $q$ and exponent $y+1$.
No extra request is needed for this argument.

\begin{proposition}[Acyclic sign recovery]\label{prop:signs}
Suppose the signed power interface is available, $t\ge8$, $F\ge10^8$,
and the unchanged raw digit and endpoint equations hold with
\eqref{eq:affine}. Then the system forces
\[
 C,A_i\ge0,\quad H\ge2,\quad B\ne0,\quad x>0,\quad |q|\ge3,
 \quad y\ge F,
\]
and
\begin{equation}\label{eq:growth}
 x\ge F+1,\qquad x\ge2+(t+1)y,\qquad
 |q|=|BH|^x\ge2^x,\qquad |B|<|q|/2.
\end{equation}
\end{proposition}
\begin{proof}
The numerator requests at base $q+1$ and Lemma~\ref{lem:power} first
give nonnegative actual CRT representatives. Hence $H\ge2$. The raw
index row $Bh_B=C+F+1$ excludes $B=0$, so $|BH|\ge2$. Its request
gives $x\ge0$ and the exact equality \eqref{eq:radix}.

Before knowing the sign of $y$, its odd multiplier gives $|y|\ge F$.
Thus $|1+(t+1)y|>1$, and \eqref{eq:last} rules out $x=0$.
Consequently $x>0$ and $|q|\ge2$. Since $q\equiv1\pmod4$, in fact
$|q|\ge3$. The retained request $q^{y+1}$ now gives $y+1\ge0$.
An odd nonzero multiple of $F>1$ cannot lie in $[-1,0]$, so $y\ge F$.

If $u_0\le-1$, \eqref{eq:first} would imply $x\le F-y\le0$.
Therefore $u_0\ge0$ and $x\ge F+1$. Equation~\eqref{eq:last} then
gives $u_t\ge1$ and the second lower bound in \eqref{eq:growth}.
Finally $|BH|\ge2$, $x\ge3$, and $H\ge2$ give the remaining bounds.
\end{proof}

Within this particular unconditional numerator-base argument, modulus
four is minimal. A progression of modulus one, two, or three cannot
avoid all three consecutive forbidden numerator bases $-1,0,1$.
The progression $q+1\equiv2\pmod4$ does avoid them. This is a small,
precisely scoped optimality observation, not a lower bound on general
Diophantine representations.

\section{A smaller frame and an affine proof horizon}

The sparser comparison profile retains $t=16+h^2$ and
$F=[2(t+1)]^t$. The main profiles instead use
\begin{equation}\label{eq:horizons}
             t=8h,\qquad d=t+2,\qquad F=d^t,
             \qquad h\in\Z.
\end{equation}
Three compatible changes keep the resulting coefficients small: bound
all affine atoms by one coordinate frame, remove a repeated horizon
term from the orientation error pool, and use the smaller growth base
$t+2$. The proof below establishes orientation before using any sign
of the horizon.

\subsection{A coordinate frame for affine atoms}

\begin{lemma}[Integer affine majorants]\label{lem:frame}
Let $L(\boldsymbol u)=c+\sum_i a_i u_i$ have integer coefficients.
Define
\[
 b(L)=|c|+\sum_i\left\lfloor a_i^2/4\right\rfloor.
\]
Then, for all integer $\boldsymbol u$,
\[
 |L(\boldsymbol u)|\le\sum_{a_i\ne0}u_i^2+b(L).
\]
\end{lemma}
\begin{proof}
For integer $a,u$, both $au$ and $-au$ are at most
$u^2+\lfloor a^2/4\rfloor$, by maximizing $au-u^2$ over the integers.
Use the triangle inequality and sum these bounds after collecting
coefficients of the same coordinate.
\end{proof}

Resolve only the new affine source identities for $t$, $d$, $R$, and $q$.
The complete source has 37 distinct requested affine atoms, using 33
underlying coordinates. The supplied \texttt{frame\_certificate.json}
lists every atom, every substitution, and its exact $b(L)$. Their maximum
is 18, attained by $d=8h+2$. We replace the old frame by
\begin{equation}\label{eq:coordinateframe}
                   W=19+\sum_{u\in\mathcal U}u^2,
\end{equation}
where $\mathcal U$ is precisely that set of 33 coordinates. Therefore
$W\ge19$ and every actual base, exponent, and output atom has absolute
value strictly less than $W$. The frame coordinates are semantic inputs
and output registers, all chosen before the private Pell witnesses;
there is no dependence on $A$ or a Pell ordinate.

For examples, the frame costs of $t=8h$, $t+1$, $F+t$, $q=4z+1$,
$q+1$, and $d=8h+2$ are respectively $16,17,16,5,6,18$.
An ordinary register $u$ costs zero; a translated register $u+1$ costs one.
The lemma covers all 37 atoms uniformly. Unlike the former frame,
\eqref{eq:coordinateframe} does not assert that each affine atom's
\emph{square} is below $W$; the following argument needs only the
established absolute-value bounds.

\subsection{An invertible change of orientation errors}

Let $e_t=Y_t^{\rm proposed}-t$ and
$e_{F+t}=Y_{F+t}^{\rm proposed}-(F+t)$ denote two of the nominal
ordinate errors. Keep $e_t$ and replace the second by
\begin{equation}\label{eq:errordifference}
              e_{F+t}-e_t
                 =Y_{F+t}^{\rm proposed}-Y_t^{\rm proposed}-F.
\end{equation}
Keep every other ordinate error and every abscissa error. Squaring and
pooling this new list enforces exactly the same zero-error conditions,
because the change of error basis is invertible over $\Z$. It removes
one repeated copy of the horizon from the quadratic pool.

Here is the unbounded soundness estimate; positive definiteness alone
would not suffice to authenticate a congruence. The retained signed
anchor and its cap first bound every classified index by the absolute
nominal anchor index, hence strictly below $W$. Each ordinary small
ordinate error has absolute value less than $2W$, and each abscissa
error has absolute value at most two. In the new difference term,
the two classified signed indices and the separate atom $F$ each
have absolute value below $W$, so its absolute value is less than
$3W$. With at most fifteen pairs, the reduced nonnegative error obeys
\[
       E<65W^2+60<4W^3<|A-1|\qquad(W\ge19).
\]
The orientation row makes $A-1$ divide $E$, so $E=0$. Invertibility
of \eqref{eq:errordifference} recovers each original ordinate error
as zero as well, and both coordinates of every Pell pair are oriented.

The exceptional anchor branch $|C_{\rm Pell}|=1$ is handled before
asserting the anchor-index identification. Its squared cap forces
nominal anchor zero and every other principal ordinate zero. Its
classified anchor index is one, and every other index is zero, so
the same size bound applies. Its anchor error is positive and is
unchanged by \eqref{eq:errordifference}; therefore $E>0$, a contradiction.

The odd frame norm still classifies at center $W+1$ with a positive
index divisible by $W$. It gives
$|A|\ge X_W(W+1)>W^W$ and $|A-1|>4W^3$. The power-modulus estimate
$|2An-n^2-1|>W^W+W$ for $|n|\ge2$ follows as before. Hence the full
signed power interface is available before the horizon has a sign.
For constructive completeness, canonical errors and their differences
are divisible by $A-1$, so the changed pooled row has its integer
quotient. The semantic-anchor cap and its dominance construction are
unchanged.

\subsection{Recovering the horizon and the required growth}

In \eqref{eq:horizons}, $d=8h+2$ has absolute value at least two for
all integer $h$. Lemma~\ref{lem:power} forces $t\ge0$ and gives
$F=d^t$. If $t=0$, then $F=1$ and $y=2v+1$ is odd. Subtracting the
two beta endpoints gives
\[
                    1=(y+1)(u_t-u_0),
\]
which is impossible because the right side is even. No sign of $y$
and no division by an endpoint modulus is used here. Thus $t$ is a
positive multiple of eight and $t\ge8$.

Now $d=t+2$ is positive and even, and
\[
 F\ge10^8,
 \qquad F>(t+1)^t,
 \qquad v_2(F)\ge t.
\]
The selector inequality $F>L_t^2$ follows from
Corollary~\ref{cor:selector} whenever $t\ge16$. The only additional
positive multiple of eight is $t=8$, for which
\[
       L_8=840,\qquad L_8^2=705\,600<10^8=F.
\]
Thus the first-digit and CRT proofs apply at every possible solution.
Proposition~\ref{prop:signs} gives
$x\ge2+(t+1)F\ge900\,000\,002>2^{29}$. This is already more than
sufficient for the local-rule majorant used below; the older
$2^{32}$ threshold was unnecessarily strong.

Finally, a single explicitly certified integral residual-row subtraction
cancels the duplicated leading horizon square in the growth-modulus
and pooled-error rows. This last change preserves their common zero
set by an elementary inverse. It is an algebraic coefficient reduction,
not another semantic assumption. It restores the maximum residual
coefficient to 64 and the collected quartic height to 7,968.

\section{Authenticating the first digit with a signed radix}

At the stage of Proposition~\ref{prop:signs}, $B$ and $q$ could still
be negative. Write $Q=|q|$ and
$c_j=\binom{F+t}{j}$. For $0\le j\le t$,
\[
 c_j\le(F^2+F)^t\le2^{t(F-1)}.
\]
The last elementary inequality holds for $F\ge7$. Combining it with
$Q\ge2^x$ and $x\ge2+(t+1)F$ gives
\[
       \frac{\max_{j\le t}c_j}{Q-1}<\frac1{32},
       \qquad c_t<Q/32.
\]

The retained common denominator is $M=q^x$. Each digit denominator
$D_i=q^{k_i}$ divides $M$, and the source uses its integer quotient
$v_i$ together with normalized remainder $\rho_i$. Its one shared bound is
\begin{equation}\label{eq:norm}
 M^2=1+4\sum_{i=1}^{11}\rho_i^2
          +\alpha^2+\beta^2+\gamma(\gamma+1).
\end{equation}
All final coordinates are integers, so the last three terms are
nonnegative. Thus $M\ne0$ and $|\rho_i|<|M|/2$.

The scaled digit row implies $v_i\mid\rho_i$: every other term is
divisible by $v_i$, since $M=v_iD_i$. Accordingly
$r_i=\rho_i/v_i$ is an \emph{integer}, and
$|r_i|<|D_i|/2$. This integrality point prevents a rational-relaxation
gap in the normalized-remainder argument.

For the first digit the signed low polynomial is
\[
                    r_* =\sum_{j=0}^{t-1}c_jq^j.
\]
Its absolute value is less than $Q^t/32$, and it is congruent to
$(q+1)^{F+t}$ modulo $q^t$. The certified remainder has the same
congruence and absolute value less than $Q^t/2$. Their difference has
absolute value less than $Q^t$, so it is zero. The exact digit row now
gives $B\equiv c_t\pmod q$. Since
\[
             |B-c_t|<Q/2+Q/32<Q,
\]
we obtain $B=c_t=\binom{F+t}{t}>0$. Only now does
$q=(BH)^x$ imply $q>0$. No earlier step assumed a positive radix.

\section{From the remaining digits to a finite refutation}

Once $B,q>0$, the predecessor's remaining range proof applies. For an
ordinary nonnegative upper representative $a$ and $0\le j\le k\le x$,
\[
             \binom aj\le\max(1,a)^k\le H^x.
\]
This includes zero representatives. The retained derived-label
majorant has exponent $\ell=2^{22}$. The endpoint bound and $F\ge10^8>2^{22}$
give $\ell(t+1)<x$, so that label's required coefficients obey the same
bound. Since $q=B^xH^x$, all eleven true normalized remainders are
less than $M/32$ and are recovered uniquely by \eqref{eq:norm}.

For the local numerical rule polynomial the retained majorant is
$\LL<\mathcal R^E$ with
$\mathcal R=2+t+x+y+F$ and $E=2^{23}$. The bounds
$y\ge F$, $x\ge2+(t+1)y$, and $t\ge8$ imply $\mathcal R<2x$.
For integers $x\ge2^{29}$,
\[
                     (2x)^E<2^x.
\]
At $x=2^{29}$, its logarithmic inequality is
$30\cdot2^{23}<2^{29}$. One then observes
$(1+1/x)^E<1/(1-E/x)<2$ to induct. Our endpoint bound gives
$x\ge900\,000\,002>2^{29}$. Thus every required local value is
strictly below the prime-power precision $2^x$.

The selector facts following Corollary~\ref{cor:selector} now recover
every bounded row. To see the essential argument, a selector
$p^e\Vert B$ with $p^e>t$ identifies its row by the raw congruence
$B\mid C+F+1$. The binomial range conditions recover all bounded
columns modulo $p^{xe}$. The local polynomial is divisible by
$p^{xe}\ge2^x$ and lies in $[0,2^x)$, so it is zero.

The algebraic boundary $y=F$ is still possible in the affine beta
formula and must be excluded. The first endpoint would give initial
accumulator zero modulo $y+1$, since $F+1=y+1$. Every local update has
the form
\[
 S_{i+1}(1+v_iF)=S_i(1+p_iF)(1+q_iF),
\]
with positive label weights. Zero propagates through all rows, contrary
to the final accumulator one. Hence $y>F$, and the initial accumulator
is the intended $F+1$.

Multiplying the updates gives an equality of products of weights
$1+jF$. The conclusion side has at most $t$ nonunit factors; the other
side, including the initial target, has at most $2t+1$. If the counts
differ, their smaller value $r$ is at most $t$. The larger product is
at least $(F+1)^{r+1}$ and the smaller at most
$t^r(F+1)^r$. This contradicts $F+1>t^t$. Thus both counts are equal
and at most $t$. Replace $F$ by an indeterminate. Every coefficient
sum is at most $(t+1)^t<F$. A nonzero difference polynomial cannot
vanish at $F$: its lowest nonzero coefficient would have to be
divisible by $F$ despite having smaller absolute value. Unique
factorization therefore recovers the label multiplicities.

For the single-conclusion flow step, take the closure of truth under
the finitely many certified rule instances. If the target were outside
this closure, each rule with conclusion outside would consume at least
one premise outside. The total number of outside conclusions minus
outside premises would be nonpositive, whereas the target balance
makes it one. This contradiction yields a finite C222 derivation.
The named logical calibration completes
the soundness direction of Theorem~\ref{thm:main}.

\section{Constructive completeness}

Start with a finite calibrated C222 refutation. Pad with truth steps to
a sufficiently large permitted horizon: $16+h^2$ for the square-horizon
profile, or a positive multiple of 8 for the affine profiles.
The horizon is chosen beyond the proof size and its finite local
parameters. Define $F$ by the corresponding growth request.

Construct the pending-multiset accumulators using the weights $1+jF$.
Let $O$ be the odd part of $t!$. Choose a sufficiently large positive
odd multiple $w$ of $O$ and set
\[
                y=Fw,\qquad v=(w-1)/2.
\]
Every selected growth base is even, so $v_2(F)\ge t>v_2(t!)$.
Consequently $y$ is an even multiple of $t!$. It can be chosen larger
than every accumulator and larger than $F$. The beta moduli
$1+(i+1)y$, $0\le i\le t$, are pairwise coprime and odd. CRT supplies
a code $x$; adding their odd product permits an arbitrarily large
\emph{even} code while preserving all residues.

Set $B=\binom{F+t}{t}$ and $N=B^x$. For $1\le i\le t$,
$v_2(F)>v_2(i)$, so $v_2(F+i)=v_2(i)$. The numerator and denominator
in the binomial therefore have equal 2-adic valuation, and $B$ is odd.
The generalized binomial CRT theorem supplies all representative
classes modulo $N^2$.

Choose the distinguished schema representative $a$ odd and all other
actual representatives even. The odd period $N^2$ permits these choices.
Choose $a$ in $[2N^2,4N^2)$ and the others positive and at most $2N^2$.
The unchanged constructor identities make the semantic anchor even
and at least $2a^2$, so it exceeds twice every other required index.
The front index $F+t$ is at most $B$ for $t\ge2$, and the other front
indices are smaller than $N^2$. The bounds used in the retained anchor
construction therefore remain valid.

Now $H=2+C+\sum A_i$ is odd. Since $B,H$ are odd and $x$ is even,
$q=(BH)^x$ is an odd square; in particular $q\equiv1\pmod4$.
Thus $z=(q-1)/4$ is an integer. This uses precisely the predecessor's
parity freedom; no new modulus-three condition is needed.

Choose the canonical shared Pell witnesses and normalized remainders.
All eleven true remainders satisfy $0\le\rho_i<M/32$, so
$4\sum\rho_i^2<M^2$. The nonnegative integer gap in \eqref{eq:norm}
is representable as $\alpha^2+\beta^2+\gamma(\gamma+1)$.
For example, Legendre's three-square theorem applied to $4g+1$ gives
this representation of any $g\ge0$: exactly one square is odd modulo
four. All unchanged private certificates are supplied by the retained
constructive Pell argument. This constructs an integer zero of the
new raw source, and the exact algebraic projections give zeros of the
delivered quartics.

Notice that $F$ is chosen before the large beta multiplier $w$.
Identifying $y$ with $F+1$ would destroy this freedom. An accumulator
with two pending nonunit labels already has value at least $(F+1)^2$,
which cannot fit below $y=F+1$. That tempting further simplification
was rejected for this precise completeness obstruction.

\section{From the source certificate to the literal quartics}

The semantic changes preserve \emph{existence} by the two refutation
translations above. They are not asserted to give a pointwise bijection
with every zero of the predecessor polynomial. In contrast, the later
finite algebraic steps give exact existential projections with specified
integral restorations; invertible row and coordinate changes give
bijections. They neither add nor lose the existence of integer zeros.

The supplied producer records the raw edits, inherited Horn
factorizations, monic coordinate eliminations, integral translations,
affine residual exposures, and reversible residual row operations.
Every eliminated coordinate has an explicit integer-polynomial
restoration. Each elementary row operation is invertible over the
integers. Finally a sum of real squares vanishes exactly when every
quadratic residual vanishes.

The final sparse serialization uses a term $[c,[i,j,\ldots]]$ for
$c x_i x_j\cdots$. The empty index list denotes a constant; repeated
indices denote repeated factors. Thus these files specify actual
ordinary integer polynomials independently of any informal semantic
register name.

\subsection{Verification scope and reproduction}

The frozen finite checks passed in ordinary and optimized Python modes.
The ordinary run also regenerated all 37 generated JSON files in a fresh
temporary directory and compared them byte for byte with the archive.
No external Python package is required for the polynomial checks.

\begin{itemize}[leftmargin=*]
\item The source checker verifies every new source edit, the complete
      37-atom/33-coordinate frame ledger, 18 inherited Horn identities,
      and 205 or 206 monic projections, depending on the profile.
\item Each final endpoint has six coordinate shears, five affine-exposure
      operations, and three or four affine projections. Every final
      residual row operation is checked exactly over the integers.
\item The independent collector recomputed all six delivered quartics
      (two intermediate systems and four final endpoints), including
      every collected coefficient and every claimed metric.
\item The source/map checker rejected 16 deliberate corruptions across
      its two profiles. The separate polynomial collector rejected
      another 18 corrupted residual, coefficient, and metric examples.
\item The literal LaTeX appendices and plain-text residual lists are
      regenerated from the frozen equations; the export check requires
      exact equality with the supplied text.
\end{itemize}

From the unpacked package root, run:
\begin{verbatim}
python3 -B verify.py
python3 -B verify.py --rebuild
python3 -B -O code/verify_endpoints.py
sh build_report.sh
\end{verbatim}
The first command checks the archive manifest and all finite certificates.
The second additionally rebuilds the 37 JSON files. The third verifies
that the checks do not depend on Python assertions. The last command
requires a standard pdfLaTeX installation and rebuilds this report.


The independent root collector imports no producer or inherited
compiler. It recomputes every coefficient of every delivered sum of
squares, checks all reported dimensions, degrees, heights and supports,
and requires corrupted coefficient, residual, and metric examples to
fail. A separate checker validates the raw source changes and the
complete finite algebraic transformations. The executable receipts
state what was actually run.

These exact finite checks establish that the delivered data implement
the specified polynomial transformations. The unbounded Pell,
prime-power, CRT, sign, and flow arguments above establish the
mathematical meaning. Neither kind of evidence is presented as a
substitute for the other. The full transitive Lean/Rocq development was
not rebuilt, and this packet does not contain a newly checked encoded
object-language ZFC derivation.

\section{Further research questions}

\begin{enumerate}[leftmargin=*]
\item Can the small square-frame estimate support further affine
      parameters in the private Pell certificates? Each proposal must
      exclude zero or unit bases before invoking signed exponent
      recovery, and must retain the anchor's squared divisibility.
\item Can the common remainder norm reuse an existing universal ternary
      representation without restricting its integer gap? Replacing a
      three-variable nonnegativity certificate by a two-square demand
      would need a new completeness theorem; size alone is insufficient.
\item Can a different finite proof calculus reduce the maximal pending
      multiset width and thereby lower the beta scale? The failed
      $y=F+1$ attempt identifies the concrete obstruction that such a
      calculus would need to remove.
\item Can exact lattice changes reduce the support or lower the coefficient
      height below 7,968 while retaining the affine horizon's smaller dimension?
      Higher-degree ideal consequences or
      different source coordinates would be a substantive next step.
\item Can the new raw source and its sign-recovery order be added to the
      repository's paused proof-assistant development as a small
      independent theorem? The best first target is the combined
      bootstrap, before attempting the full polynomial compilation.
\end{enumerate}

\section*{Attribution and integration}
The new work is the compatible affine bootstrap and linear selector
composition, its proofs, literal endpoints, and finite certificates.
The numerical calculus, logical calibration, factorial-free CRT
architecture, shared Pell frame, normalized remainder construction,
cofinal parity choices, and earlier affine and row reductions are
inherited and credited. The package retains the predecessor and its
source notices. Independent review here means separate AI-assisted
mathematical and algebraic reviews, not external human refereeing.
No upstream branch, issue, pull request, or repository file was modified.

\begin{thebibliography}{9}
\bibitem{repo} V. Reshetnikov, \emph{ProveIt}, pinned repository snapshot
\repo. \url{https://github.com/VladimirReshetnikov/ProveIt}.
\bibitem{parent} ProveIt incoming research packet,
\emph{Removing factorial authenticity from Diophantine consistency},
25 October 2026. Retained verbatim in \texttt{inherited\_factorial\_free/}.
\bibitem{shared} ProveIt incoming research packet,
\emph{Shared factorial errors: explicit Diophantine equations for Con(ZFC)},
25 October 2026. Comparison identity recorded in \texttt{PROVENANCE.json}.
\bibitem{farhi} B. Farhi, \emph{An identity involving the least common
multiple of binomial coefficients and its application}, American
Mathematical Monthly 116 (2009), 836--839.
\url{https://arxiv.org/abs/0906.2295}.
\bibitem{chan} H.-L. Chan and M. Norrish,
\emph{Proof Pearl: Bounding Least Common Multiples with Triangles},
Journal of Automated Reasoning 62 (2019), 171--192.
\url{https://doi.org/10.1007/s10817-017-9438-0}.
\bibitem{bayer} J. Bayer, M. David, M. Ha\ss ler, Y. Matiyasevich and
D. Schleicher, \emph{Diophantine Equations over $\Z$: Universal Bounds
and Parallel Formalization}, arXiv:2506.20909, version 3,
22 October 2026. \url{https://arxiv.org/abs/2506.20909v3}.
\end{thebibliography}

\clearpage\appendix
\section{The complete main residual list}\label{app:residuals}
Every displayed coordinate is an unrestricted integer. This appendix
defines all residuals of \eqref{eq:main}; no omitted family or implicit
coefficient parameter is required. The sparser comparison list is supplied
in its data directory.
\begingroup\footnotesize
\setlength{\jot}{2pt}\allowdisplaybreaks[4]
\begin{align*}
R_{0} &={} -x_{1} -8x_{1}x_{212} +x_{6} +8x_{212} \\
R_{1} &={} -x_{1}x_{4} -x_{2} +x_{5}x_{6} \\
R_{2} &={} -1 +x_{0} -x_{5}x_{6} \\
R_{3} &={} 1 +x_{2} -x_{3}x_{17} +x_{15} \\
R_{4} &={} -1 -x_{7}^{2} -2x_{7}x_{8} -x_{8}^{2} +x_{18} \\
R_{5} &={} -10x_{7} +x_{7}^{2} -6x_{7}x_{8} -8x_{7}x_{18} -20x_{8} \\
&\quad{} -11x_{8}^{2} -16x_{8}x_{18} -15x_{18} -4x_{18}^{2} +x_{19} \\
R_{6} &={} -2 +2x_{7} -x_{7}^{2} -4x_{7}x_{19} -4x_{19} \\
&\quad{} -4x_{19}^{2} +x_{20} \\
R_{7} &={} 1 +6x_{7} -36x_{7}^{2} -24x_{7}x_{20} +2x_{20} \\
&\quad{} -4x_{20}^{2} +x_{21} \\
R_{8} &={} -2 -x_{8} -x_{8}^{2} -2x_{8}x_{9} -x_{9} \\
&\quad{} -x_{9}^{2} +x_{22} \\
R_{9} &={} -8 +7x_{7} -x_{7}^{2} +4x_{7}x_{9} -4x_{7}x_{22} \\
&\quad{} -4x_{9} +4x_{22} +x_{23} -x_{45} \\
R_{10} &={} -1 +x_{7} -x_{7}^{2} -2x_{7}x_{9} +x_{9} \\
&\quad{} -x_{9}^{2} +x_{24} \\
R_{11} &={} 23 +16x_{7} -8x_{7}x_{8} -8x_{7}x_{9} -16x_{7}x_{18} \\
&\quad{} -8x_{7}x_{24} +19x_{8} +9x_{8}^{2} -36x_{8}x_{9} -12x_{8}x_{24} \\
&\quad{} +16x_{9} +8x_{9}^{2} -24x_{9}x_{18} -8x_{18}x_{24} -x_{19} \\
&\quad{} +x_{25} -x_{30} \\
R_{12} &={} -4x_{23}^{2} -8x_{23}x_{25} -4x_{25}^{2} +x_{26} \\
R_{13} &={} -2 -8x_{7} +8x_{7}x_{22} +8x_{7}x_{24} -8x_{8} \\
&\quad{} +8x_{8}x_{22} +8x_{8}x_{24} -4x_{18} +8x_{22} -4x_{22}^{2} \\
&\quad{} -8x_{22}x_{24} +8x_{24} -4x_{24}^{2} +x_{27} \\
R_{14} &={} 6 +4x_{7} +32x_{7}x_{8} +16x_{7}x_{18} +16x_{7}x_{24} \\
&\quad{} +8x_{7}x_{27} +24x_{8} -32x_{8}x_{24} -16x_{8}x_{27} +12x_{18} \\
&\quad{} -16x_{18}x_{24} -8x_{18}x_{27} -x_{19} +8x_{24} -16x_{24}^{2} \\
&\quad{} -16x_{24}x_{27} +4x_{27} -4x_{27}^{2} +x_{28} \\
R_{15} &={} 22 -5x_{7} -x_{7}^{2} -8x_{7}x_{8} -12x_{7}x_{9} \\
&\quad{} -4x_{7}x_{22} -8x_{8} -16x_{8}x_{9} -16x_{8}x_{22} -20x_{9} \\
&\quad{} -20x_{9}^{2} -24x_{9}x_{22} -28x_{22} -4x_{22}^{2} +x_{29} \\
R_{16} &={} 19 -25x_{7} -7x_{7}^{2} -8x_{7}x_{8} -30x_{7}x_{9} \\
&\quad{} -16x_{7}x_{24} +3x_{8} +5x_{8}^{2} -4x_{8}x_{24} -27x_{9} \\
&\quad{} -21x_{9}^{2} -24x_{9}x_{24} -6x_{22} -17x_{24} -4x_{24}^{2} \\
&\quad{} +x_{30} \\
R_{17} &={} 2 -4x_{29}^{2} -8x_{29}x_{30} -4x_{30}^{2} +x_{31} \\
R_{18} &={} -19x_{7} +8x_{7}^{2} +8x_{7}x_{18} +8x_{7}x_{22} -46x_{8} \\
&\quad{} -24x_{8}^{2} +24x_{8}x_{18} +24x_{8}x_{22} -2x_{9} +4x_{18} \\
&\quad{} -4x_{18}^{2} -8x_{18}x_{22} +14x_{22} -4x_{22}^{2} -x_{24} \\
&\quad{} +x_{32} \\
R_{19} &={} -3 -6x_{7} -4x_{9} -2x_{24} -4x_{32} \\
&\quad{} -4x_{32}^{2} +x_{33} \\
R_{20} &={} 8 +42x_{7} +7x_{7}^{2} +28x_{7}x_{19} +48x_{8} \\
&\quad{} -32x_{8}^{2} -32x_{8}x_{18} +64x_{8}x_{19} -40x_{18}^{2} +32x_{18}x_{19} \\
&\quad{} +14x_{19} -4x_{19}^{2} +x_{34} +2x_{37} \\
R_{21} &={} 1 +6x_{7} -36x_{7}^{2} -24x_{7}x_{34} +2x_{34} \\
&\quad{} -4x_{34}^{2} +x_{35} \\
R_{22} &={} -x_{8}^{2} -4x_{8}x_{35} -4x_{35}^{2} +x_{36} \\
R_{23} &={} -6 -2x_{7} -4x_{8} -14x_{18} -16x_{18}^{2} \\
&\quad{} +x_{37} \\
R_{24} &={} 24 +21x_{8} -21x_{8}^{2} -18x_{8}x_{9} -24x_{8}x_{22} \\
&\quad{} -24x_{8}x_{37} +3x_{9} -9x_{9}^{2} -16x_{9}x_{37} +3x_{22} \\
&\quad{} -12x_{22}^{2} -8x_{22}x_{37} +4x_{37} -4x_{37}^{2} +x_{38} \\
&\quad{} +2x_{45} \\
R_{25} &={} 16 -10x_{7} -20x_{7}^{2} -16x_{7}x_{9} -24x_{7}x_{24} \\
&\quad{} -24x_{7}x_{38} -12x_{9} -16x_{9}x_{24} -16x_{9}x_{38} -14x_{24} \\
&\quad{} -4x_{24}^{2} -8x_{24}x_{38} +4x_{38} -4x_{38}^{2} +x_{39} \\
R_{26} &={} -x_{9}^{2} -2x_{9}x_{10} -x_{10}^{2} +x_{40} \\
R_{27} &={} 2 +x_{7}^{2} -4x_{7}x_{40} +2x_{9} -2x_{9}^{2} \\
&\quad{} -8x_{9}x_{40} -2x_{24} +4x_{40} -4x_{40}^{2} +x_{41} \\
R_{28} &={} -31 -16x_{9} -16x_{9}x_{22} +16x_{9}x_{24} -8x_{22} \\
&\quad{} +12x_{22}^{2} -8x_{22}x_{24} +24x_{24} -4x_{24}^{2} +x_{42} \\
&\quad{} -4x_{45} \\
R_{29} &={} -4x_{9} -16x_{24}^{2} -32x_{24}x_{40} -16x_{24}x_{41} -16x_{24}x_{42} \\
&\quad{} -4x_{40} -16x_{40}^{2} -16x_{40}x_{41} -16x_{40}x_{42} -2x_{41} \\
&\quad{} -4x_{41}^{2} -8x_{41}x_{42} -4x_{42}^{2} +x_{43} \\
R_{30} &={} -16x_{7}^{2} +8x_{7}x_{9} +32x_{7}x_{24} +16x_{7}x_{27} +x_{9} \\
&\quad{} -x_{9}^{2} -8x_{9}x_{24} -4x_{9}x_{27} -16x_{24}^{2} -16x_{24}x_{27} \\
&\quad{} -4x_{27}^{2} +x_{44} \\
R_{31} &={} -1 -8x_{9} -4x_{9}^{2} +8x_{9}x_{22} +8x_{22} \\
&\quad{} -4x_{22}^{2} +x_{45} \\
R_{32} &={} -1 -9x_{9}^{2} +24x_{9}x_{22} +12x_{9}x_{45} -16x_{22}^{2} \\
&\quad{} -16x_{22}x_{45} -4x_{45}^{2} +x_{46} \\
R_{33} &={} 1 +x_{8} -x_{8}^{2} -4x_{8}x_{9} -4x_{8}x_{46} \\
&\quad{} -4x_{9}^{2} -8x_{9}x_{46} -4x_{46}^{2} +x_{47} \\
R_{34} &={} -1 -4x_{47} -4x_{47}^{2} +x_{48} \\
R_{35} &={} -1 -x_{9}^{2} -8x_{9}x_{47} -4x_{9}x_{48} -16x_{47}^{2} \\
&\quad{} -16x_{47}x_{48} -4x_{48}^{2} +x_{49} \\
R_{36} &={} 2 +x_{8} -x_{8}^{2} -4x_{8}x_{9} -4x_{8}x_{49} \\
&\quad{} -4x_{9}^{2} -8x_{9}x_{49} -4x_{49}^{2} +x_{50} \\
R_{37} &={} -4x_{44}^{2} -8x_{44}x_{50} -4x_{50}^{2} +x_{51} \\
R_{38} &={} -9 +7x_{7} -x_{7}^{2} -8x_{7}x_{44} -4x_{7}x_{51} \\
&\quad{} +24x_{44} -16x_{44}^{2} -16x_{44}x_{51} +12x_{51} -4x_{51}^{2} \\
&\quad{} +x_{52} \\
R_{39} &={} -1 -4x_{52} -4x_{52}^{2} +x_{53} \\
R_{40} &={} -x_{9}^{2} -8x_{9}x_{52} -4x_{9}x_{53} -16x_{52}^{2} -16x_{52}x_{53} \\
&\quad{} -4x_{53}^{2} +x_{54} \\
R_{41} &={} -x_{7} -x_{7}^{2} -4x_{7}x_{9} -4x_{7}x_{54} -4x_{9} \\
&\quad{} -4x_{9}^{2} -8x_{9}x_{54} -4x_{54} -4x_{54}^{2} +x_{55} \\
R_{42} &={} -10 -12x_{55} -4x_{55}^{2} +x_{56} \\
R_{43} &={} 2 +4x_{9} +16x_{9}x_{55} +8x_{9}x_{56} -4x_{22} \\
&\quad{} -16x_{22}x_{55} -8x_{22}x_{56} -x_{45} +8x_{55} -16x_{55}^{2} \\
&\quad{} -16x_{55}x_{56} +4x_{56} -4x_{56}^{2} +x_{57} \\
R_{44} &={} x_{8} -x_{8}^{2} +8x_{8}x_{9} -8x_{8}x_{22} -4x_{8}x_{57} \\
&\quad{} -16x_{9}^{2} +32x_{9}x_{22} +16x_{9}x_{57} -16x_{22}^{2} -16x_{22}x_{57} \\
&\quad{} -4x_{57}^{2} +x_{58} \\
R_{45} &={} -10 -12x_{50} -4x_{50}^{2} +x_{59} \\
R_{46} &={} 1 -16x_{50}^{2} -16x_{50}x_{58} -16x_{50}x_{59} -4x_{58}^{2} \\
&\quad{} -8x_{58}x_{59} -4x_{59}^{2} +x_{60} \\
R_{47} &={} x_{9} -x_{9}^{2} -8x_{9}x_{58} -4x_{9}x_{60} -16x_{58}^{2} \\
&\quad{} -16x_{58}x_{60} -4x_{60}^{2} +x_{61} \\
R_{48} &={} -6x_{7} -3x_{7}^{2} -18x_{7}x_{8} -16x_{8} -27x_{8}^{2} \\
&\quad{} +4x_{8}x_{9} +3x_{9} +3x_{9}^{2} -4x_{9}x_{18} -19x_{18} \\
&\quad{} -8x_{18}^{2} +x_{19} +x_{62} -4x_{75} \\
R_{49} &={} 1 -8x_{7} -4x_{7}x_{9} +8x_{7}x_{62} -16x_{8} \\
&\quad{} -8x_{8}x_{9} +16x_{8}x_{62} +x_{9} +x_{9}^{2} +4x_{9}x_{18} \\
&\quad{} +8x_{18} -8x_{18}x_{62} +3x_{62} -4x_{62}^{2} +x_{63} \\
R_{50} &={} -2 -12x_{7} -16x_{7}x_{9} +16x_{7}x_{61} +16x_{7}x_{63} \\
&\quad{} -24x_{8} -32x_{8}x_{9} +32x_{8}x_{61} +32x_{8}x_{63} +4x_{9} \\
&\quad{} +4x_{9}^{2} +16x_{9}x_{18} +12x_{18} -16x_{18}x_{61} -16x_{18}x_{63} \\
&\quad{} -4x_{61}^{2} -8x_{61}x_{63} -4x_{62} -2x_{63} -4x_{63}^{2} \\
&\quad{} +x_{64} \\
R_{51} &={} -2 -4x_{64} -4x_{64}^{2} +x_{65} \\
R_{52} &={} -x_{2}x_{12} +x_{66} \\
R_{53} &={} x_{11} -x_{12} -2x_{21}x_{66} -x_{67} \\
R_{54} &={} x_{11} -x_{12} -4x_{23}x_{66} -2x_{26}x_{66} -x_{68} \\
R_{55} &={} -2x_{28}x_{66} +2x_{43}x_{66} -2x_{66} -x_{69} +x_{74} \\
R_{56} &={} x_{11} -x_{12} -4x_{29}x_{66} -2x_{31}x_{66} -x_{70} \\
R_{57} &={} x_{11} -x_{12} -2x_{33}x_{66} +6x_{66} -x_{71} \\
R_{58} &={} x_{11} -x_{12} -4x_{35}x_{66} -2x_{36}x_{66} -x_{72} \\
R_{59} &={} x_{11} -x_{12} -2x_{39}x_{66} -x_{73} \\
R_{60} &={} x_{11} -x_{12} -2x_{43}x_{66} -4x_{66} -x_{74} \\
R_{61} &={} -x_{7}x_{8} +x_{7}x_{9} +x_{8} +x_{8}x_{9} -x_{9} \\
&\quad{} -x_{9}^{2} +x_{75} \\
R_{62} &={} 1 +2x_{10} +x_{10}^{2} -x_{75}^{2} +x_{76} \\
R_{63} &={} -x_{76}^{2} +x_{77} \\
R_{64} &={} -4x_{64}x_{66} -2x_{65}x_{66} +2x_{66} +x_{78} \\
R_{65} &={} -x_{11}^{2} +2x_{11}x_{12} +2x_{11}x_{78} -x_{12}^{2} -2x_{12}x_{78} \\
&\quad{} -x_{78}^{2} +x_{79} \\
R_{66} &={} -1 -x_{2}x_{8} +x_{80} \\
R_{67} &={} -x_{12}x_{80} +x_{82}x_{83} +x_{84} \\
R_{68} &={} -1 +x_{2} -x_{2}x_{7} +x_{81} \\
R_{69} &={} -1 +2x_{2} +2x_{2}x_{7} +4x_{2}x_{8} -2x_{2}x_{18} \\
&\quad{} +x_{82} \\
R_{70} &={} -x_{11}x_{81} +x_{83} \\
R_{71} &={} -6x_{2}x_{12} +x_{11}x_{80} +12x_{12} -8x_{12}x_{80} -4x_{12}x_{81} \\
&\quad{} -x_{12}x_{82} +x_{85} \\
R_{72} &={} -x_{11}x_{67} +x_{12}x_{67} +x_{86} \\
R_{73} &={} -x_{68}x_{86} +x_{87} \\
R_{74} &={} -x_{69}x_{87} +x_{88} \\
R_{75} &={} -x_{70}x_{88} +x_{89} \\
R_{76} &={} -x_{71}x_{89} +x_{90} \\
R_{77} &={} -x_{72}x_{90} +x_{91} \\
R_{78} &={} -x_{73}x_{91} +x_{92} \\
R_{79} &={} -x_{74}x_{92} +x_{93} \\
R_{80} &={} -x_{84}x_{93} +x_{94} \\
R_{81} &={} -x_{85}x_{94} +x_{95} \\
R_{82} &={} -x_{95}^{2} +x_{96} \\
R_{83} &={} -x_{1} -x_{1}x_{15} +x_{15} +x_{97} \\
R_{84} &={} -x_{13}x_{97} +x_{98} \\
R_{85} &={} -2x_{0}^{2} +2x_{0}x_{11} +2x_{0}x_{12} +2x_{0}x_{98} +2x_{0}x_{99} \\
&\quad{} -x_{11}^{2} -2x_{11}x_{98} -x_{12}^{2} -2x_{12}x_{99} +x_{16}x_{101} \\
&\quad{} -x_{77}x_{96} -x_{79}x_{96} -x_{98}^{2} -x_{99}^{2} \\
R_{86} &={} -x_{1}x_{14} +x_{14} -x_{14}x_{97} +x_{99} \\
R_{87} &={} -x_{7} -x_{8} -4x_{18} -2x_{21} -4x_{23} \\
&\quad{} -2x_{26} -2x_{28} -4x_{29} -2x_{31} -2x_{33} \\
&\quad{} -4x_{35} -2x_{36} -2x_{39} -2x_{43} -4x_{64} \\
&\quad{} -2x_{65} +x_{100} \\
R_{88} &={} -x_{104} -4x_{104}x_{236} +x_{106} \\
R_{89} &={} -x_{3}x_{128} +x_{102}x_{227} -x_{105} -x_{106}x_{128} \\
R_{90} &={} x_{107}x_{228} -x_{109}x_{131} -x_{110} \\
R_{91} &={} x_{111}x_{228} -x_{112}x_{131} -x_{113} \\
R_{92} &={} x_{114}x_{228} -x_{115}x_{131} -x_{116} \\
R_{93} &={} x_{117}x_{228} -x_{118}x_{131} -x_{119} \\
R_{94} &={} x_{120}x_{229} -x_{122}x_{131} -x_{123} \\
R_{95} &={} x_{124}x_{229} -x_{125}x_{131} -x_{126} \\
R_{96} &={} -x_{16}x_{128} +x_{131} \\
R_{97} &={} x_{127} -x_{129}x_{131} -x_{130} \\
R_{98} &={} -x_{131}x_{133} +x_{132} -x_{134} \\
R_{99} &={} -x_{131}x_{136} +x_{135}x_{227} -x_{137} \\
R_{100} &={} -x_{131}x_{139} +x_{138}x_{228} -x_{140} \\
R_{101} &={} x_{141} -2x_{141}x_{143} +x_{144} \\
R_{102} &={} -x_{141}x_{144} -x_{144} +x_{145} \\
R_{103} &={} x_{144}^{2} -x_{145}^{2} +x_{165} \\
R_{104} &={} -x_{142}x_{147} -x_{147} +x_{148} \\
R_{105} &={} 2x_{146} +x_{146}^{2} +x_{147}^{2} -x_{148}^{2} \\
R_{106} &={} 1 -4x_{142} -16x_{142}x_{212} +x_{149} +16x_{212} \\
&\quad{} +64x_{212}^{2} \\
R_{107} &={} -1 +x_{2} +x_{142}x_{147} -x_{146} -x_{147} \\
&\quad{} -8x_{147}x_{212} +x_{149}x_{150} \\
R_{108} &={} 1 -2x_{3} +x_{3}^{2} -2x_{3}x_{142} +x_{151} \\
R_{109} &={} -x_{3}x_{196} +x_{16} +x_{151}x_{152} -x_{195} +x_{197} \\
R_{110} &={} -1 +2x_{142}x_{211} -x_{153} +2x_{211} -x_{211}^{2} \\
R_{111} &={} 1 +x_{153}x_{154} -x_{195} -x_{196}x_{211} +x_{197} \\
&\quad{} +4x_{236} \\
R_{112} &={} -x_{142}x_{156} -x_{156} +x_{157} \\
R_{113} &={} 2x_{155} +x_{155}^{2} +x_{156}^{2} -x_{157}^{2} \\
R_{114} &={} 1 -4x_{142} -8x_{142}x_{236} +x_{158} +8x_{236} \\
&\quad{} +16x_{236}^{2} \\
R_{115} &={} -1 +x_{102} -x_{155} -2x_{156} -4x_{156}x_{236} \\
&\quad{} +x_{157} +x_{158}x_{159} \\
R_{116} &={} -1 +x_{103} -2x_{142}x_{160} -x_{146} -x_{147} \\
&\quad{} -4x_{147}x_{236} +x_{148} +x_{158}x_{160} +x_{160} +8x_{160}x_{236} \\
R_{117} &={} 1 +x_{142} -x_{142}x_{162} -x_{162} +x_{163} \\
R_{118} &={} 1 +2x_{161} +x_{161}^{2} -2x_{162} +x_{162}^{2} \\
&\quad{} -x_{163}^{2} \\
R_{119} &={} 1 +x_{107} +x_{158}x_{164} -x_{161} -2x_{162} \\
&\quad{} -4x_{162}x_{236} +x_{163} +4x_{236} \\
R_{120} &={} -2x_{142} -x_{142}^{2} +x_{165} \\
R_{121} &={} -1 +x_{108} +x_{142}x_{148} -2x_{142}x_{166} -x_{146} \\
&\quad{} -4x_{146}x_{236} -x_{147}x_{165} -4x_{148}x_{236} +x_{158}x_{166} +x_{166} \\
&\quad{} +8x_{166}x_{236} -4x_{236} \\
R_{122} &={} -x_{142}x_{168} -x_{168} +x_{169} \\
R_{123} &={} 2x_{167} +x_{167}^{2} +x_{168}^{2} -x_{169}^{2} \\
R_{124} &={} -1 +x_{111} +x_{158}x_{170} -x_{167} -2x_{168} \\
&\quad{} -4x_{168}x_{236} +x_{169} \\
R_{125} &={} -x_{142}x_{172} -x_{172} +x_{173} \\
R_{126} &={} 2x_{171} +x_{171}^{2} +x_{172}^{2} -x_{173}^{2} \\
R_{127} &={} -1 +x_{114} +x_{158}x_{174} -x_{171} -2x_{172} \\
&\quad{} -4x_{172}x_{236} +x_{173} \\
R_{128} &={} -x_{142}x_{176} -x_{176} +x_{177} \\
R_{129} &={} 2x_{175} +x_{175}^{2} +x_{176}^{2} -x_{177}^{2} \\
R_{130} &={} -1 +x_{117} +x_{158}x_{178} -x_{175} -2x_{176} \\
&\quad{} -4x_{176}x_{236} +x_{177} \\
R_{131} &={} -x_{142}x_{180} -x_{180} +x_{181} \\
R_{132} &={} 2x_{179} +x_{179}^{2} +x_{180}^{2} -x_{181}^{2} \\
R_{133} &={} -1 +x_{120} +x_{158}x_{182} -x_{179} -2x_{180} \\
&\quad{} -4x_{180}x_{236} +x_{181} \\
R_{134} &={} -x_{142}x_{184} -x_{184} +x_{185} \\
R_{135} &={} 2x_{183} +x_{183}^{2} +x_{184}^{2} -x_{185}^{2} \\
R_{136} &={} -1 +x_{121} -2x_{142}x_{186} +x_{158}x_{186} -x_{183} \\
&\quad{} -x_{184} -4x_{184}x_{236} +x_{185} +x_{186} +8x_{186}x_{236} \\
R_{137} &={} -x_{142}x_{188} -x_{188} +x_{189} \\
R_{138} &={} 2x_{187} +x_{187}^{2} +x_{188}^{2} -x_{189}^{2} \\
R_{139} &={} -1 +x_{124} +x_{158}x_{190} -x_{187} -2x_{188} \\
&\quad{} -4x_{188}x_{236} +x_{189} \\
R_{140} &={} -x_{142}x_{192} -x_{192} +x_{193} \\
R_{141} &={} 2x_{191} +x_{191}^{2} +x_{192}^{2} -x_{193}^{2} \\
R_{142} &={} -1 +x_{127} +x_{158}x_{194} -x_{191} -2x_{192} \\
&\quad{} -4x_{192}x_{236} +x_{193} \\
R_{143} &={} 1 -x_{142}x_{196} -x_{196} +x_{197} \\
R_{144} &={} -1 +2x_{195} +x_{195}^{2} +x_{196}^{2} -2x_{197} \\
&\quad{} -x_{197}^{2} \\
R_{145} &={} x_{128} -2x_{142}x_{198} +x_{158}x_{198} -x_{195} -x_{196} \\
&\quad{} -4x_{196}x_{236} +x_{197} +x_{198} +8x_{198}x_{236} \\
R_{146} &={} -x_{142}x_{200} -x_{200} +x_{201} \\
R_{147} &={} 2x_{199} +x_{199}^{2} +x_{200}^{2} -x_{201}^{2} \\
R_{148} &={} -1 +x_{132} +x_{158}x_{202} -x_{199} -2x_{200} \\
&\quad{} -4x_{200}x_{236} +x_{201} \\
R_{149} &={} -x_{142}x_{204} -x_{204} +x_{205} \\
R_{150} &={} 2x_{203} +x_{203}^{2} +x_{204}^{2} -x_{205}^{2} \\
R_{151} &={} -1 +x_{135} +x_{158}x_{206} -x_{203} -2x_{204} \\
&\quad{} -4x_{204}x_{236} +x_{205} \\
R_{152} &={} -x_{142}x_{208} -x_{208} +x_{209} \\
R_{153} &={} 2x_{207} +x_{207}^{2} +x_{208}^{2} -x_{209}^{2} \\
R_{154} &={} -1 +x_{138} +x_{158}x_{210} -x_{207} -2x_{208} \\
&\quad{} -4x_{208}x_{236} +x_{209} \\
R_{155} &={} -32 -x_{0}^{2} +2x_{1} -x_{1}^{2} -x_{2}^{2} \\
&\quad{} -x_{3}^{2} +4x_{7} +2x_{7}x_{8} +2x_{8} +2x_{9} \\
&\quad{} -x_{9}^{2} +2x_{10} -x_{10}^{2} +2x_{11} -x_{11}^{2} \\
&\quad{} +2x_{12} -x_{12}^{2} +2x_{13} -x_{13}^{2} +2x_{14} \\
&\quad{} -x_{14}^{2} +2x_{15} -x_{15}^{2} -x_{16}^{2} -x_{18} \\
&\quad{} +2x_{100} -x_{100}^{2} -x_{102}^{2} -x_{103}^{2} -x_{107}^{2} \\
&\quad{} -x_{108}^{2} -x_{111}^{2} -x_{114}^{2} -x_{117}^{2} -x_{120}^{2} \\
&\quad{} -x_{121}^{2} -x_{124}^{2} -x_{127}^{2} -x_{128}^{2} -x_{132}^{2} \\
&\quad{} -x_{135}^{2} -x_{138}^{2} +x_{141} -x_{211}^{2} -x_{212}^{2} \\
&\quad{} -x_{236}^{2} \\
R_{156} &={} -x_{3} -x_{3}x_{7} -x_{3}x_{8} -x_{3}x_{9} -x_{3}x_{10} \\
&\quad{} -x_{3}x_{11} -x_{3}x_{12} -x_{3}x_{13} -x_{3}x_{14} -x_{3}x_{15} \\
&\quad{} +x_{211} \\
R_{157} &={} -x_{142}x_{214} -x_{214} +x_{223} \\
R_{158} &={} -1 +x_{213}^{2} +x_{214}^{2} -x_{223}^{2} \\
R_{159} &={} -x_{216} -x_{216}x_{217} +x_{224} \\
R_{160} &={} -1 +x_{215}^{2} +x_{216}^{2} -x_{224}^{2} \\
R_{161} &={} -x_{208}^{2} +x_{225} \\
R_{162} &={} x_{214} -2x_{219}x_{225} -x_{225} \\
R_{163} &={} -4x_{214} -8x_{214}x_{218} +x_{217} \\
R_{164} &={} -x_{142} -x_{213}x_{220} +x_{217} \\
R_{165} &={} -1 -x_{207} -x_{213}x_{221} +x_{215} \\
R_{166} &={} -x_{100} -4x_{208}x_{222} +x_{216} \\
R_{167} &={} -3 +x_{0}^{2} -2x_{0}x_{196} +x_{1}^{2} -2x_{1}x_{184} \\
&\quad{} +x_{2}^{2} +2x_{2}x_{147} -2x_{2}x_{156} +x_{7}^{2} -2x_{7}x_{162} \\
&\quad{} +x_{8}^{2} -2x_{8}x_{168} +x_{9}^{2} -2x_{9}x_{172} +x_{10}^{2} \\
&\quad{} -2x_{10}x_{176} +x_{11}^{2} -2x_{11}x_{180} +x_{12}^{2} -2x_{12}x_{188} \\
&\quad{} +x_{13}^{2} -2x_{13}x_{192} +x_{14}^{2} -2x_{14}x_{200} +x_{15}^{2} \\
&\quad{} -2x_{15}x_{204} -2x_{100}x_{208} +4x_{142} +16x_{142}x_{212} -x_{142}x_{226} \\
&\quad{} +x_{146}^{2} +x_{147}^{2} -2x_{147}x_{156} -16x_{147}x_{212} -x_{149} \\
&\quad{} +x_{155}^{2} +x_{161}^{2} +2x_{162} +x_{167}^{2} +x_{171}^{2} \\
&\quad{} +x_{175}^{2} +x_{179}^{2} +x_{183}^{2} +x_{187}^{2} +x_{191}^{2} \\
&\quad{} +x_{195}^{2} +x_{199}^{2} +x_{203}^{2} +x_{207}^{2} +2x_{208}^{2} \\
&\quad{} -16x_{212} -x_{230}^{2} -x_{231}^{2} -x_{232} -x_{232}^{2} \\
R_{168} &={} -x_{103}x_{227} +x_{128} \\
R_{169} &={} -x_{108}x_{228} +x_{128} \\
R_{170} &={} -x_{121}x_{229} +x_{128} \\
R_{171} &={} -1 +x_{1} -x_{2} -2x_{2}x_{237} \\
R_{172} &={} -2 -x_{100}^{2} -x_{147}^{2} -x_{156}^{2} +2x_{162} \\
&\quad{} -x_{162}^{2} -x_{168}^{2} -x_{172}^{2} -x_{176}^{2} -x_{180}^{2} \\
&\quad{} -x_{184}^{2} -x_{188}^{2} -x_{192}^{2} -x_{196}^{2} -x_{200}^{2} \\
&\quad{} -x_{204}^{2} +x_{208}^{2} -x_{230}^{2} -x_{231}^{2} -x_{232} \\
&\quad{} -x_{232}^{2} \\
R_{173} &={} -1 -4x_{105}^{2} -4x_{110}^{2} -4x_{113}^{2} -4x_{116}^{2} \\
&\quad{} -4x_{119}^{2} -4x_{123}^{2} -4x_{126}^{2} +x_{128}^{2} -4x_{130}^{2} \\
&\quad{} -4x_{134}^{2} -4x_{137}^{2} -4x_{140}^{2} -x_{233}^{2} -x_{234}^{2} \\
&\quad{} -x_{235} -x_{235}^{2} \\
\end{align*}
\endgroup

\clearpage
\section{The complete \SmallVars-variable residual list}\label{app:smallest}
For the endpoint named \texttt{minimum\_dimension}, define
\[
 P_{\SmallVars}(x_0,\ldots,x_{\SmallLastVar})
 =\sum_{j=0}^{\SmallLastRow}S_j(x_0,\ldots,x_{\SmallLastVar})^2.
\]
The coordinates below belong to this endpoint's own integral coordinate
system. They are independent of the coordinates in the preceding appendix.
All \SmallRows\ residuals are printed in full.
\begingroup\footnotesize
\setlength{\jot}{2pt}\allowdisplaybreaks[4]
\begin{align*}
S_{0} &={} -x_{1} -8x_{1}x_{211} +x_{6} +8x_{211} \\
S_{1} &={} -x_{1}x_{4} -x_{2} +x_{5}x_{6} \\
S_{2} &={} -1 +x_{0} -x_{5}x_{6} \\
S_{3} &={} 1 +x_{2} -x_{3}x_{17} +x_{15} \\
S_{4} &={} -1 -x_{7}^{2} -2x_{7}x_{8} -x_{8}^{2} +x_{18} \\
S_{5} &={} -10x_{7} +x_{7}^{2} -6x_{7}x_{8} -8x_{7}x_{18} -20x_{8} \\
&\quad{} -11x_{8}^{2} -16x_{8}x_{18} -15x_{18} -4x_{18}^{2} +x_{19} \\
S_{6} &={} -2 +2x_{7} -x_{7}^{2} -4x_{7}x_{19} -4x_{19} \\
&\quad{} -4x_{19}^{2} +x_{20} \\
S_{7} &={} 1 +6x_{7} -36x_{7}^{2} -24x_{7}x_{20} +2x_{20} \\
&\quad{} -4x_{20}^{2} +x_{21} \\
S_{8} &={} -2 -x_{8} -x_{8}^{2} -2x_{8}x_{9} -x_{9} \\
&\quad{} -x_{9}^{2} +x_{22} \\
S_{9} &={} -8 +7x_{7} -x_{7}^{2} +4x_{7}x_{9} -4x_{7}x_{22} \\
&\quad{} -4x_{9} +4x_{22} +x_{23} -x_{45} \\
S_{10} &={} -1 +x_{7} -x_{7}^{2} -2x_{7}x_{9} +x_{9} \\
&\quad{} -x_{9}^{2} +x_{24} \\
S_{11} &={} 23 +16x_{7} -8x_{7}x_{8} -8x_{7}x_{9} -16x_{7}x_{18} \\
&\quad{} -8x_{7}x_{24} +19x_{8} +9x_{8}^{2} -36x_{8}x_{9} -12x_{8}x_{24} \\
&\quad{} +16x_{9} +8x_{9}^{2} -24x_{9}x_{18} -8x_{18}x_{24} -x_{19} \\
&\quad{} +x_{25} -x_{30} \\
S_{12} &={} -4x_{23}^{2} -8x_{23}x_{25} -4x_{25}^{2} +x_{26} \\
S_{13} &={} -2 -8x_{7} +8x_{7}x_{22} +8x_{7}x_{24} -8x_{8} \\
&\quad{} +8x_{8}x_{22} +8x_{8}x_{24} -4x_{18} +8x_{22} -4x_{22}^{2} \\
&\quad{} -8x_{22}x_{24} +8x_{24} -4x_{24}^{2} +x_{27} \\
S_{14} &={} 6 +4x_{7} +32x_{7}x_{8} +16x_{7}x_{18} +16x_{7}x_{24} \\
&\quad{} +8x_{7}x_{27} +24x_{8} -32x_{8}x_{24} -16x_{8}x_{27} +12x_{18} \\
&\quad{} -16x_{18}x_{24} -8x_{18}x_{27} -x_{19} +8x_{24} -16x_{24}^{2} \\
&\quad{} -16x_{24}x_{27} +4x_{27} -4x_{27}^{2} +x_{28} \\
S_{15} &={} 22 -5x_{7} -x_{7}^{2} -8x_{7}x_{8} -12x_{7}x_{9} \\
&\quad{} -4x_{7}x_{22} -8x_{8} -16x_{8}x_{9} -16x_{8}x_{22} -20x_{9} \\
&\quad{} -20x_{9}^{2} -24x_{9}x_{22} -28x_{22} -4x_{22}^{2} +x_{29} \\
S_{16} &={} 19 -25x_{7} -7x_{7}^{2} -8x_{7}x_{8} -30x_{7}x_{9} \\
&\quad{} -16x_{7}x_{24} +3x_{8} +5x_{8}^{2} -4x_{8}x_{24} -27x_{9} \\
&\quad{} -21x_{9}^{2} -24x_{9}x_{24} -6x_{22} -17x_{24} -4x_{24}^{2} \\
&\quad{} +x_{30} \\
S_{17} &={} 2 -4x_{29}^{2} -8x_{29}x_{30} -4x_{30}^{2} +x_{31} \\
S_{18} &={} -19x_{7} +8x_{7}^{2} +8x_{7}x_{18} +8x_{7}x_{22} -46x_{8} \\
&\quad{} -24x_{8}^{2} +24x_{8}x_{18} +24x_{8}x_{22} -2x_{9} +4x_{18} \\
&\quad{} -4x_{18}^{2} -8x_{18}x_{22} +14x_{22} -4x_{22}^{2} -x_{24} \\
&\quad{} +x_{32} \\
S_{19} &={} -3 -6x_{7} -4x_{9} -2x_{24} -4x_{32} \\
&\quad{} -4x_{32}^{2} +x_{33} \\
S_{20} &={} 8 +42x_{7} +7x_{7}^{2} +28x_{7}x_{19} +48x_{8} \\
&\quad{} -32x_{8}^{2} -32x_{8}x_{18} +64x_{8}x_{19} -40x_{18}^{2} +32x_{18}x_{19} \\
&\quad{} +14x_{19} -4x_{19}^{2} +x_{34} +2x_{37} \\
S_{21} &={} 1 +6x_{7} -36x_{7}^{2} -24x_{7}x_{34} +2x_{34} \\
&\quad{} -4x_{34}^{2} +x_{35} \\
S_{22} &={} -x_{8}^{2} -4x_{8}x_{35} -4x_{35}^{2} +x_{36} \\
S_{23} &={} -6 -2x_{7} -4x_{8} -14x_{18} -16x_{18}^{2} \\
&\quad{} +x_{37} \\
S_{24} &={} 24 +21x_{8} -21x_{8}^{2} -18x_{8}x_{9} -24x_{8}x_{22} \\
&\quad{} -24x_{8}x_{37} +3x_{9} -9x_{9}^{2} -16x_{9}x_{37} +3x_{22} \\
&\quad{} -12x_{22}^{2} -8x_{22}x_{37} +4x_{37} -4x_{37}^{2} +x_{38} \\
&\quad{} +2x_{45} \\
S_{25} &={} 16 -10x_{7} -20x_{7}^{2} -16x_{7}x_{9} -24x_{7}x_{24} \\
&\quad{} -24x_{7}x_{38} -12x_{9} -16x_{9}x_{24} -16x_{9}x_{38} -14x_{24} \\
&\quad{} -4x_{24}^{2} -8x_{24}x_{38} +4x_{38} -4x_{38}^{2} +x_{39} \\
S_{26} &={} -x_{9}^{2} -2x_{9}x_{10} -x_{10}^{2} +x_{40} \\
S_{27} &={} 2 +x_{7}^{2} -4x_{7}x_{40} +2x_{9} -2x_{9}^{2} \\
&\quad{} -8x_{9}x_{40} -2x_{24} +4x_{40} -4x_{40}^{2} +x_{41} \\
S_{28} &={} -31 -16x_{9} -16x_{9}x_{22} +16x_{9}x_{24} -8x_{22} \\
&\quad{} +12x_{22}^{2} -8x_{22}x_{24} +24x_{24} -4x_{24}^{2} +x_{42} \\
&\quad{} -4x_{45} \\
S_{29} &={} -4x_{9} -16x_{24}^{2} -32x_{24}x_{40} -16x_{24}x_{41} -16x_{24}x_{42} \\
&\quad{} -4x_{40} -16x_{40}^{2} -16x_{40}x_{41} -16x_{40}x_{42} -2x_{41} \\
&\quad{} -4x_{41}^{2} -8x_{41}x_{42} -4x_{42}^{2} +x_{43} \\
S_{30} &={} -16x_{7}^{2} +8x_{7}x_{9} +32x_{7}x_{24} +16x_{7}x_{27} +x_{9} \\
&\quad{} -x_{9}^{2} -8x_{9}x_{24} -4x_{9}x_{27} -16x_{24}^{2} -16x_{24}x_{27} \\
&\quad{} -4x_{27}^{2} +x_{44} \\
S_{31} &={} -1 -8x_{9} -4x_{9}^{2} +8x_{9}x_{22} +8x_{22} \\
&\quad{} -4x_{22}^{2} +x_{45} \\
S_{32} &={} -1 -9x_{9}^{2} +24x_{9}x_{22} +12x_{9}x_{45} -16x_{22}^{2} \\
&\quad{} -16x_{22}x_{45} -4x_{45}^{2} +x_{46} \\
S_{33} &={} 1 +x_{8} -x_{8}^{2} -4x_{8}x_{9} -4x_{8}x_{46} \\
&\quad{} -4x_{9}^{2} -8x_{9}x_{46} -4x_{46}^{2} +x_{47} \\
S_{34} &={} -1 -4x_{47} -4x_{47}^{2} +x_{48} \\
S_{35} &={} -1 -x_{9}^{2} -8x_{9}x_{47} -4x_{9}x_{48} -16x_{47}^{2} \\
&\quad{} -16x_{47}x_{48} -4x_{48}^{2} +x_{49} \\
S_{36} &={} 2 +x_{8} -x_{8}^{2} -4x_{8}x_{9} -4x_{8}x_{49} \\
&\quad{} -4x_{9}^{2} -8x_{9}x_{49} -4x_{49}^{2} +x_{50} \\
S_{37} &={} -4x_{44}^{2} -8x_{44}x_{50} -4x_{50}^{2} +x_{51} \\
S_{38} &={} -9 +7x_{7} -x_{7}^{2} -8x_{7}x_{44} -4x_{7}x_{51} \\
&\quad{} +24x_{44} -16x_{44}^{2} -16x_{44}x_{51} +12x_{51} -4x_{51}^{2} \\
&\quad{} +x_{52} \\
S_{39} &={} -1 -4x_{52} -4x_{52}^{2} +x_{53} \\
S_{40} &={} -x_{9}^{2} -8x_{9}x_{52} -4x_{9}x_{53} -16x_{52}^{2} -16x_{52}x_{53} \\
&\quad{} -4x_{53}^{2} +x_{54} \\
S_{41} &={} -x_{7} -x_{7}^{2} -4x_{7}x_{9} -4x_{7}x_{54} -4x_{9} \\
&\quad{} -4x_{9}^{2} -8x_{9}x_{54} -4x_{54} -4x_{54}^{2} +x_{55} \\
S_{42} &={} -10 -12x_{55} -4x_{55}^{2} +x_{56} \\
S_{43} &={} 2 +4x_{9} +16x_{9}x_{55} +8x_{9}x_{56} -4x_{22} \\
&\quad{} -16x_{22}x_{55} -8x_{22}x_{56} -x_{45} +8x_{55} -16x_{55}^{2} \\
&\quad{} -16x_{55}x_{56} +4x_{56} -4x_{56}^{2} +x_{57} \\
S_{44} &={} x_{8} -x_{8}^{2} +8x_{8}x_{9} -8x_{8}x_{22} -4x_{8}x_{57} \\
&\quad{} -16x_{9}^{2} +32x_{9}x_{22} +16x_{9}x_{57} -16x_{22}^{2} -16x_{22}x_{57} \\
&\quad{} -4x_{57}^{2} +x_{58} \\
S_{45} &={} -10 -12x_{50} -4x_{50}^{2} +x_{59} \\
S_{46} &={} 1 -16x_{50}^{2} -16x_{50}x_{58} -16x_{50}x_{59} -4x_{58}^{2} \\
&\quad{} -8x_{58}x_{59} -4x_{59}^{2} +x_{60} \\
S_{47} &={} x_{9} -x_{9}^{2} -8x_{9}x_{58} -4x_{9}x_{60} -16x_{58}^{2} \\
&\quad{} -16x_{58}x_{60} -4x_{60}^{2} +x_{61} \\
S_{48} &={} -6x_{7} -3x_{7}^{2} -18x_{7}x_{8} -16x_{8} -27x_{8}^{2} \\
&\quad{} +4x_{8}x_{9} +3x_{9} +3x_{9}^{2} -4x_{9}x_{18} -19x_{18} \\
&\quad{} -8x_{18}^{2} +x_{19} +x_{62} -4x_{75} \\
S_{49} &={} 1 -8x_{7} -4x_{7}x_{9} +8x_{7}x_{62} -16x_{8} \\
&\quad{} -8x_{8}x_{9} +16x_{8}x_{62} +x_{9} +x_{9}^{2} +4x_{9}x_{18} \\
&\quad{} +8x_{18} -8x_{18}x_{62} +3x_{62} -4x_{62}^{2} +x_{63} \\
S_{50} &={} -2 -12x_{7} -16x_{7}x_{9} +16x_{7}x_{61} +16x_{7}x_{63} \\
&\quad{} -24x_{8} -32x_{8}x_{9} +32x_{8}x_{61} +32x_{8}x_{63} +4x_{9} \\
&\quad{} +4x_{9}^{2} +16x_{9}x_{18} +12x_{18} -16x_{18}x_{61} -16x_{18}x_{63} \\
&\quad{} -4x_{61}^{2} -8x_{61}x_{63} -4x_{62} -2x_{63} -4x_{63}^{2} \\
&\quad{} +x_{64} \\
S_{51} &={} -2 -4x_{64} -4x_{64}^{2} +x_{65} \\
S_{52} &={} -x_{2}x_{12} +x_{66} \\
S_{53} &={} x_{11} -x_{12} -2x_{21}x_{66} -x_{67} \\
S_{54} &={} x_{11} -x_{12} -4x_{23}x_{66} -2x_{26}x_{66} -x_{68} \\
S_{55} &={} -2x_{28}x_{66} +2x_{43}x_{66} -2x_{66} -x_{69} +x_{74} \\
S_{56} &={} x_{11} -x_{12} -4x_{29}x_{66} -2x_{31}x_{66} -x_{70} \\
S_{57} &={} x_{11} -x_{12} -2x_{33}x_{66} +6x_{66} -x_{71} \\
S_{58} &={} x_{11} -x_{12} -4x_{35}x_{66} -2x_{36}x_{66} -x_{72} \\
S_{59} &={} x_{11} -x_{12} -2x_{39}x_{66} -x_{73} \\
S_{60} &={} x_{11} -x_{12} -2x_{43}x_{66} -4x_{66} -x_{74} \\
S_{61} &={} -x_{7}x_{8} +x_{7}x_{9} +x_{8} +x_{8}x_{9} -x_{9} \\
&\quad{} -x_{9}^{2} +x_{75} \\
S_{62} &={} 1 +2x_{10} +x_{10}^{2} -x_{75}^{2} +x_{76} \\
S_{63} &={} -x_{76}^{2} +x_{77} \\
S_{64} &={} -4x_{64}x_{66} -2x_{65}x_{66} +2x_{66} +x_{78} \\
S_{65} &={} -x_{11}^{2} +2x_{11}x_{12} +2x_{11}x_{78} -x_{12}^{2} -2x_{12}x_{78} \\
&\quad{} -x_{78}^{2} +x_{79} \\
S_{66} &={} -1 -x_{2}x_{8} +x_{80} \\
S_{67} &={} -x_{12}x_{80} +x_{82}x_{83} +x_{84} \\
S_{68} &={} -1 +x_{2} -x_{2}x_{7} +x_{81} \\
S_{69} &={} -1 +2x_{2} +2x_{2}x_{7} +4x_{2}x_{8} -2x_{2}x_{18} \\
&\quad{} +x_{82} \\
S_{70} &={} -x_{11}x_{81} +x_{83} \\
S_{71} &={} -6x_{2}x_{12} +x_{11}x_{80} +12x_{12} -8x_{12}x_{80} -4x_{12}x_{81} \\
&\quad{} -x_{12}x_{82} +x_{85} \\
S_{72} &={} -x_{11}x_{67} +x_{12}x_{67} +x_{86} \\
S_{73} &={} -x_{68}x_{86} +x_{87} \\
S_{74} &={} -x_{69}x_{87} +x_{88} \\
S_{75} &={} -x_{70}x_{88} +x_{89} \\
S_{76} &={} -x_{71}x_{89} +x_{90} \\
S_{77} &={} -x_{72}x_{90} +x_{91} \\
S_{78} &={} -x_{73}x_{91} +x_{92} \\
S_{79} &={} -x_{74}x_{92} +x_{93} \\
S_{80} &={} -x_{84}x_{93} +x_{94} \\
S_{81} &={} -x_{85}x_{94} +x_{95} \\
S_{82} &={} -x_{95}^{2} +x_{96} \\
S_{83} &={} -x_{1} -x_{1}x_{15} +x_{15} +x_{97} \\
S_{84} &={} -x_{13}x_{97} +x_{98} \\
S_{85} &={} -2x_{0}^{2} +2x_{0}x_{11} +2x_{0}x_{12} +2x_{0}x_{98} +2x_{0}x_{99} \\
&\quad{} -x_{11}^{2} -2x_{11}x_{98} -x_{12}^{2} -2x_{12}x_{99} +x_{16}x_{100} \\
&\quad{} -x_{77}x_{96} -x_{79}x_{96} -x_{98}^{2} -x_{99}^{2} \\
S_{86} &={} -x_{1}x_{14} +x_{14} -x_{14}x_{97} +x_{99} \\
S_{87} &={} -x_{103} -4x_{103}x_{235} +x_{105} \\
S_{88} &={} -x_{3}x_{127} +x_{101}x_{226} -x_{104} -x_{105}x_{127} \\
S_{89} &={} x_{106}x_{227} -x_{108}x_{130} -x_{109} \\
S_{90} &={} x_{110}x_{227} -x_{111}x_{130} -x_{112} \\
S_{91} &={} x_{113}x_{227} -x_{114}x_{130} -x_{115} \\
S_{92} &={} x_{116}x_{227} -x_{117}x_{130} -x_{118} \\
S_{93} &={} x_{119}x_{228} -x_{121}x_{130} -x_{122} \\
S_{94} &={} x_{123}x_{228} -x_{124}x_{130} -x_{125} \\
S_{95} &={} -x_{16}x_{127} +x_{130} \\
S_{96} &={} x_{126} -x_{128}x_{130} -x_{129} \\
S_{97} &={} -x_{130}x_{132} +x_{131} -x_{133} \\
S_{98} &={} -x_{130}x_{135} +x_{134}x_{226} -x_{136} \\
S_{99} &={} -x_{130}x_{138} +x_{137}x_{227} -x_{139} \\
S_{100} &={} x_{140} -2x_{140}x_{142} +x_{143} \\
S_{101} &={} -x_{140}x_{143} -x_{143} +x_{144} \\
S_{102} &={} x_{143}^{2} -x_{144}^{2} +x_{164} \\
S_{103} &={} -x_{141}x_{146} -x_{146} +x_{147} \\
S_{104} &={} 2x_{145} +x_{145}^{2} +x_{146}^{2} -x_{147}^{2} \\
S_{105} &={} 1 -4x_{141} -16x_{141}x_{211} +x_{148} +16x_{211} \\
&\quad{} +64x_{211}^{2} \\
S_{106} &={} -1 +x_{2} +x_{141}x_{146} -x_{145} -x_{146} \\
&\quad{} -8x_{146}x_{211} +x_{148}x_{149} \\
S_{107} &={} 1 -2x_{3} +x_{3}^{2} -2x_{3}x_{141} +x_{150} \\
S_{108} &={} -x_{3}x_{195} +x_{16} +x_{150}x_{151} -x_{194} +x_{196} \\
S_{109} &={} -1 +2x_{141}x_{210} -x_{152} +2x_{210} -x_{210}^{2} \\
S_{110} &={} 1 +x_{152}x_{153} -x_{194} -x_{195}x_{210} +x_{196} \\
&\quad{} +4x_{235} \\
S_{111} &={} -x_{141}x_{155} -x_{155} +x_{156} \\
S_{112} &={} 2x_{154} +x_{154}^{2} +x_{155}^{2} -x_{156}^{2} \\
S_{113} &={} 1 -4x_{141} -8x_{141}x_{235} +x_{157} +8x_{235} \\
&\quad{} +16x_{235}^{2} \\
S_{114} &={} -1 +x_{101} -x_{154} -2x_{155} -4x_{155}x_{235} \\
&\quad{} +x_{156} +x_{157}x_{158} \\
S_{115} &={} -1 +x_{102} -2x_{141}x_{159} -x_{145} -x_{146} \\
&\quad{} -4x_{146}x_{235} +x_{147} +x_{157}x_{159} +x_{159} +8x_{159}x_{235} \\
S_{116} &={} 1 +x_{141} -x_{141}x_{161} -x_{161} +x_{162} \\
S_{117} &={} 1 +2x_{160} +x_{160}^{2} -2x_{161} +x_{161}^{2} \\
&\quad{} -x_{162}^{2} \\
S_{118} &={} 1 +x_{106} +x_{157}x_{163} -x_{160} -2x_{161} \\
&\quad{} -4x_{161}x_{235} +x_{162} +4x_{235} \\
S_{119} &={} -2x_{141} -x_{141}^{2} +x_{164} \\
S_{120} &={} -1 +x_{107} +x_{141}x_{147} -2x_{141}x_{165} -x_{145} \\
&\quad{} -4x_{145}x_{235} -x_{146}x_{164} -4x_{147}x_{235} +x_{157}x_{165} +x_{165} \\
&\quad{} +8x_{165}x_{235} -4x_{235} \\
S_{121} &={} -x_{141}x_{167} -x_{167} +x_{168} \\
S_{122} &={} 2x_{166} +x_{166}^{2} +x_{167}^{2} -x_{168}^{2} \\
S_{123} &={} -1 +x_{110} +x_{157}x_{169} -x_{166} -2x_{167} \\
&\quad{} -4x_{167}x_{235} +x_{168} \\
S_{124} &={} -x_{141}x_{171} -x_{171} +x_{172} \\
S_{125} &={} 2x_{170} +x_{170}^{2} +x_{171}^{2} -x_{172}^{2} \\
S_{126} &={} -1 +x_{113} +x_{157}x_{173} -x_{170} -2x_{171} \\
&\quad{} -4x_{171}x_{235} +x_{172} \\
S_{127} &={} -x_{141}x_{175} -x_{175} +x_{176} \\
S_{128} &={} 2x_{174} +x_{174}^{2} +x_{175}^{2} -x_{176}^{2} \\
S_{129} &={} -1 +x_{116} +x_{157}x_{177} -x_{174} -2x_{175} \\
&\quad{} -4x_{175}x_{235} +x_{176} \\
S_{130} &={} -x_{141}x_{179} -x_{179} +x_{180} \\
S_{131} &={} 2x_{178} +x_{178}^{2} +x_{179}^{2} -x_{180}^{2} \\
S_{132} &={} -1 +x_{119} +x_{157}x_{181} -x_{178} -2x_{179} \\
&\quad{} -4x_{179}x_{235} +x_{180} \\
S_{133} &={} -x_{141}x_{183} -x_{183} +x_{184} \\
S_{134} &={} 2x_{182} +x_{182}^{2} +x_{183}^{2} -x_{184}^{2} \\
S_{135} &={} -1 +x_{120} -2x_{141}x_{185} +x_{157}x_{185} -x_{182} \\
&\quad{} -x_{183} -4x_{183}x_{235} +x_{184} +x_{185} +8x_{185}x_{235} \\
S_{136} &={} -x_{141}x_{187} -x_{187} +x_{188} \\
S_{137} &={} 2x_{186} +x_{186}^{2} +x_{187}^{2} -x_{188}^{2} \\
S_{138} &={} -1 +x_{123} +x_{157}x_{189} -x_{186} -2x_{187} \\
&\quad{} -4x_{187}x_{235} +x_{188} \\
S_{139} &={} -x_{141}x_{191} -x_{191} +x_{192} \\
S_{140} &={} 2x_{190} +x_{190}^{2} +x_{191}^{2} -x_{192}^{2} \\
S_{141} &={} -1 +x_{126} +x_{157}x_{193} -x_{190} -2x_{191} \\
&\quad{} -4x_{191}x_{235} +x_{192} \\
S_{142} &={} 1 -x_{141}x_{195} -x_{195} +x_{196} \\
S_{143} &={} -1 +2x_{194} +x_{194}^{2} +x_{195}^{2} -2x_{196} \\
&\quad{} -x_{196}^{2} \\
S_{144} &={} x_{127} -2x_{141}x_{197} +x_{157}x_{197} -x_{194} -x_{195} \\
&\quad{} -4x_{195}x_{235} +x_{196} +x_{197} +8x_{197}x_{235} \\
S_{145} &={} -x_{141}x_{199} -x_{199} +x_{200} \\
S_{146} &={} 2x_{198} +x_{198}^{2} +x_{199}^{2} -x_{200}^{2} \\
S_{147} &={} -1 +x_{131} +x_{157}x_{201} -x_{198} -2x_{199} \\
&\quad{} -4x_{199}x_{235} +x_{200} \\
S_{148} &={} -x_{141}x_{203} -x_{203} +x_{204} \\
S_{149} &={} 2x_{202} +x_{202}^{2} +x_{203}^{2} -x_{204}^{2} \\
S_{150} &={} -1 +x_{134} +x_{157}x_{205} -x_{202} -2x_{203} \\
&\quad{} -4x_{203}x_{235} +x_{204} \\
S_{151} &={} -x_{141}x_{207} -x_{207} +x_{208} \\
S_{152} &={} 2x_{206} +x_{206}^{2} +x_{207}^{2} -x_{208}^{2} \\
S_{153} &={} -1 +x_{137} +x_{157}x_{209} -x_{206} -2x_{207} \\
&\quad{} -4x_{207}x_{235} +x_{208} \\
S_{154} &={} -30 -x_{0}^{2} +2x_{1} -x_{1}^{2} -x_{2}^{2} \\
&\quad{} -x_{3}^{2} +6x_{7} +2x_{7}x_{8} +4x_{8} +2x_{9} \\
&\quad{} -x_{9}^{2} +2x_{10} -x_{10}^{2} +2x_{11} -x_{11}^{2} \\
&\quad{} +2x_{12} -x_{12}^{2} +2x_{13} -x_{13}^{2} +2x_{14} \\
&\quad{} -x_{14}^{2} +2x_{15} -x_{15}^{2} -x_{16}^{2} +7x_{18} \\
&\quad{} +4x_{21} +8x_{23} +4x_{26} +4x_{28} +8x_{29} \\
&\quad{} +4x_{31} +4x_{33} +8x_{35} +4x_{36} +4x_{39} \\
&\quad{} +4x_{43} +8x_{64} +4x_{65} -x_{101}^{2} -x_{102}^{2} \\
&\quad{} -x_{106}^{2} -x_{107}^{2} -x_{110}^{2} -x_{113}^{2} -x_{116}^{2} \\
&\quad{} -x_{119}^{2} -x_{120}^{2} -x_{123}^{2} -x_{126}^{2} -x_{127}^{2} \\
&\quad{} -x_{131}^{2} -x_{134}^{2} -x_{137}^{2} +x_{140} +x_{146}^{2} \\
&\quad{} +x_{155}^{2} -2x_{161} +x_{161}^{2} +x_{167}^{2} +x_{171}^{2} \\
&\quad{} +x_{175}^{2} +x_{179}^{2} +x_{183}^{2} +x_{187}^{2} +x_{191}^{2} \\
&\quad{} +x_{195}^{2} +x_{199}^{2} +x_{203}^{2} -x_{207}^{2} -x_{210}^{2} \\
&\quad{} -x_{211}^{2} +x_{229}^{2} +x_{230}^{2} +x_{231} +x_{231}^{2} \\
&\quad{} -x_{235}^{2} \\
S_{155} &={} -x_{3} -x_{3}x_{7} -x_{3}x_{8} -x_{3}x_{9} -x_{3}x_{10} \\
&\quad{} -x_{3}x_{11} -x_{3}x_{12} -x_{3}x_{13} -x_{3}x_{14} -x_{3}x_{15} \\
&\quad{} +x_{210} \\
S_{156} &={} -x_{141}x_{213} -x_{213} +x_{222} \\
S_{157} &={} -1 +x_{212}^{2} +x_{213}^{2} -x_{222}^{2} \\
S_{158} &={} -x_{215} -x_{215}x_{216} +x_{223} \\
S_{159} &={} -1 +x_{214}^{2} +x_{215}^{2} -x_{223}^{2} \\
S_{160} &={} -x_{207}^{2} +x_{224} \\
S_{161} &={} x_{213} -2x_{218}x_{224} -x_{224} \\
S_{162} &={} -4x_{213} -8x_{213}x_{217} +x_{216} \\
S_{163} &={} -x_{141} -x_{212}x_{219} +x_{216} \\
S_{164} &={} -1 -x_{206} -x_{212}x_{220} +x_{214} \\
S_{165} &={} -x_{7} -x_{8} -4x_{18} -2x_{21} -4x_{23} \\
&\quad{} -2x_{26} -2x_{28} -4x_{29} -2x_{31} -2x_{33} \\
&\quad{} -4x_{35} -2x_{36} -2x_{39} -2x_{43} -4x_{64} \\
&\quad{} -2x_{65} -4x_{207}x_{221} +x_{215} \\
S_{166} &={} -3 +x_{0}^{2} -2x_{0}x_{195} +x_{1}^{2} -2x_{1}x_{183} \\
&\quad{} +x_{2}^{2} +2x_{2}x_{146} -2x_{2}x_{155} +x_{7}^{2} -2x_{7}x_{161} \\
&\quad{} -2x_{7}x_{207} +x_{8}^{2} -2x_{8}x_{167} -2x_{8}x_{207} +x_{9}^{2} \\
&\quad{} -2x_{9}x_{171} +x_{10}^{2} -2x_{10}x_{175} +x_{11}^{2} -2x_{11}x_{179} \\
&\quad{} +x_{12}^{2} -2x_{12}x_{187} +x_{13}^{2} -2x_{13}x_{191} +x_{14}^{2} \\
&\quad{} -2x_{14}x_{199} +x_{15}^{2} -2x_{15}x_{203} -8x_{18}x_{207} -4x_{21}x_{207} \\
&\quad{} -8x_{23}x_{207} -4x_{26}x_{207} -4x_{28}x_{207} -8x_{29}x_{207} -4x_{31}x_{207} \\
&\quad{} -4x_{33}x_{207} -8x_{35}x_{207} -4x_{36}x_{207} -4x_{39}x_{207} -4x_{43}x_{207} \\
&\quad{} -8x_{64}x_{207} -4x_{65}x_{207} +4x_{141} +16x_{141}x_{211} -x_{141}x_{225} \\
&\quad{} +x_{145}^{2} +x_{146}^{2} -2x_{146}x_{155} -16x_{146}x_{211} -x_{148} \\
&\quad{} +x_{154}^{2} +x_{160}^{2} +2x_{161} +x_{166}^{2} +x_{170}^{2} \\
&\quad{} +x_{174}^{2} +x_{178}^{2} +x_{182}^{2} +x_{186}^{2} +x_{190}^{2} \\
&\quad{} +x_{194}^{2} +x_{198}^{2} +x_{202}^{2} +x_{206}^{2} +2x_{207}^{2} \\
&\quad{} -16x_{211} -x_{229}^{2} -x_{230}^{2} -x_{231} -x_{231}^{2} \\
S_{167} &={} -x_{102}x_{226} +x_{127} \\
S_{168} &={} -x_{107}x_{227} +x_{127} \\
S_{169} &={} -x_{120}x_{228} +x_{127} \\
S_{170} &={} -1 +x_{1} -x_{2} -2x_{2}x_{236} \\
S_{171} &={} -1 -8x_{7}x_{18} -4x_{7}x_{21} -8x_{7}x_{23} -4x_{7}x_{26} \\
&\quad{} -4x_{7}x_{28} -8x_{7}x_{29} -4x_{7}x_{31} -4x_{7}x_{33} -8x_{7}x_{35} \\
&\quad{} -4x_{7}x_{36} -4x_{7}x_{39} -4x_{7}x_{43} -8x_{7}x_{64} -4x_{7}x_{65} \\
&\quad{} -8x_{8}x_{18} -4x_{8}x_{21} -8x_{8}x_{23} -4x_{8}x_{26} -4x_{8}x_{28} \\
&\quad{} -8x_{8}x_{29} -4x_{8}x_{31} -4x_{8}x_{33} -8x_{8}x_{35} -4x_{8}x_{36} \\
&\quad{} -4x_{8}x_{39} -4x_{8}x_{43} -8x_{8}x_{64} -4x_{8}x_{65} -x_{18} \\
&\quad{} -16x_{18}^{2} -16x_{18}x_{21} -32x_{18}x_{23} -16x_{18}x_{26} -16x_{18}x_{28} \\
&\quad{} -32x_{18}x_{29} -16x_{18}x_{31} -16x_{18}x_{33} -32x_{18}x_{35} -16x_{18}x_{36} \\
&\quad{} -16x_{18}x_{39} -16x_{18}x_{43} -32x_{18}x_{64} -16x_{18}x_{65} -4x_{21}^{2} \\
&\quad{} -16x_{21}x_{23} -8x_{21}x_{26} -8x_{21}x_{28} -16x_{21}x_{29} -8x_{21}x_{31} \\
&\quad{} -8x_{21}x_{33} -16x_{21}x_{35} -8x_{21}x_{36} -8x_{21}x_{39} -8x_{21}x_{43} \\
&\quad{} -16x_{21}x_{64} -8x_{21}x_{65} -16x_{23}^{2} -16x_{23}x_{26} -16x_{23}x_{28} \\
&\quad{} -32x_{23}x_{29} -16x_{23}x_{31} -16x_{23}x_{33} -32x_{23}x_{35} -16x_{23}x_{36} \\
&\quad{} -16x_{23}x_{39} -16x_{23}x_{43} -32x_{23}x_{64} -16x_{23}x_{65} -4x_{26}^{2} \\
&\quad{} -8x_{26}x_{28} -16x_{26}x_{29} -8x_{26}x_{31} -8x_{26}x_{33} -16x_{26}x_{35} \\
&\quad{} -8x_{26}x_{36} -8x_{26}x_{39} -8x_{26}x_{43} -16x_{26}x_{64} -8x_{26}x_{65} \\
&\quad{} -4x_{28}^{2} -16x_{28}x_{29} -8x_{28}x_{31} -8x_{28}x_{33} -16x_{28}x_{35} \\
&\quad{} -8x_{28}x_{36} -8x_{28}x_{39} -8x_{28}x_{43} -16x_{28}x_{64} -8x_{28}x_{65} \\
&\quad{} -16x_{29}^{2} -16x_{29}x_{31} -16x_{29}x_{33} -32x_{29}x_{35} -16x_{29}x_{36} \\
&\quad{} -16x_{29}x_{39} -16x_{29}x_{43} -32x_{29}x_{64} -16x_{29}x_{65} -4x_{31}^{2} \\
&\quad{} -8x_{31}x_{33} -16x_{31}x_{35} -8x_{31}x_{36} -8x_{31}x_{39} -8x_{31}x_{43} \\
&\quad{} -16x_{31}x_{64} -8x_{31}x_{65} -4x_{33}^{2} -16x_{33}x_{35} -8x_{33}x_{36} \\
&\quad{} -8x_{33}x_{39} -8x_{33}x_{43} -16x_{33}x_{64} -8x_{33}x_{65} -16x_{35}^{2} \\
&\quad{} -16x_{35}x_{36} -16x_{35}x_{39} -16x_{35}x_{43} -32x_{35}x_{64} -16x_{35}x_{65} \\
&\quad{} -4x_{36}^{2} -8x_{36}x_{39} -8x_{36}x_{43} -16x_{36}x_{64} -8x_{36}x_{65} \\
&\quad{} -4x_{39}^{2} -8x_{39}x_{43} -16x_{39}x_{64} -8x_{39}x_{65} -4x_{43}^{2} \\
&\quad{} -16x_{43}x_{64} -8x_{43}x_{65} -16x_{64}^{2} -16x_{64}x_{65} -4x_{65}^{2} \\
&\quad{} -x_{146}^{2} -x_{155}^{2} +2x_{161} -x_{161}^{2} -x_{167}^{2} \\
&\quad{} -x_{171}^{2} -x_{175}^{2} -x_{179}^{2} -x_{183}^{2} -x_{187}^{2} \\
&\quad{} -x_{191}^{2} -x_{195}^{2} -x_{199}^{2} -x_{203}^{2} +x_{207}^{2} \\
&\quad{} -x_{229}^{2} -x_{230}^{2} -x_{231} -x_{231}^{2} \\
S_{172} &={} -1 -4x_{104}^{2} -4x_{109}^{2} -4x_{112}^{2} -4x_{115}^{2} \\
&\quad{} -4x_{118}^{2} -4x_{122}^{2} -4x_{125}^{2} +x_{127}^{2} -4x_{129}^{2} \\
&\quad{} -4x_{133}^{2} -4x_{136}^{2} -4x_{139}^{2} -x_{232}^{2} -x_{233}^{2} \\
&\quad{} -x_{234} -x_{234}^{2} \\
\end{align*}
\endgroup

\end{document}
